% concatenated from main.tex
\documentclass[12pt]{article}%
\usepackage[applemac]{inputenc}
\usepackage{graphicx}  % Required to include images
\usepackage{caption}   % Optional: for better captions
\usepackage{color, colortbl}
\usepackage{booktabs}
\usepackage{amsmath,amssymb}
\usepackage[T1]{fontenc}
\usepackage[T1]{fontenc}
%\usepackage{showlabels}
\usepackage{amsfonts}
\usepackage[applemac]{inputenc}
\usepackage{amsmath,amssymb}
\usepackage{color}
\usepackage{color, colortbl}
\usepackage{amsmath}
\usepackage{amssymb}
\usepackage{graphicx}
\usepackage{tikz}
\usepackage{geometry}
\usetikzlibrary{arrows}
\usepackage{amsfonts}
\usepackage{amsmath,amssymb}
\usepackage[applemac]{inputenc}
\usepackage{color}
\usepackage{amsmath}
\usepackage{amssymb}
\usepackage{graphicx}
\usepackage{tikz}
\newtheorem{theorem}{Theorem}[section]
\newtheorem{theoremm}{Theorem}[subsection]
\newtheorem{lemma}[theorem]{Lemma}
\newcommand{\E}{E}
\newtheorem{proposition}[theorem]{Proposition}
\newtheorem{corollary}[theorem]{Corollary}
\newtheorem{definition}[theorem]{Definition}
\newtheorem{notation}[theorem]{Notation}
\newtheorem{example}[theorem]{Example}
\newtheorem{assumption}[theorem]{Assumption}
\newtheorem{question}[theorem]{Question}
\newtheorem{conclusion}[theorem]{Conclusion}
\newtheorem{problem}[theorem]{Problem}
%%%%%%%%%%%subsection%%%%%%%%%%%%%%%%%%
\newtheorem{propositionn}[theoremm]{Proposition}
\newtheorem{corollaryy}[theoremm]{Corollary}
\newtheorem{definitionn}[theoremm]{Definition}
\newtheorem{notationn}[theoremm]{Notation}
\newtheorem{examplee}[theoremm]{Example}
\newtheorem{assumptionn}[theoremm]{Assumption}
\newtheorem{questionn}[theoremm]{Question}
\newtheorem{conclusionn}[theoremm]{Conclusion}
\newtheorem{problemm}[theoremm]{Problem}
%\theoremstyle{definition}
\newtheorem{remark}[theorem]{Remark}
\newtheorem{remarkk}[theoremm]{Remark}


\setcounter{MaxMatrixCols}{30}
\usetikzlibrary{arrows,shapes,automata,petri}
\newenvironment{myenumerate}{\renewcommand{\theenumi}{\roman{enumi}}
\renewcommand{\labelenumi}{\bf (\theenumi)}
\begin{enumerate}}{\end{enumerate}}
\newcommand{\dproof}{\noindent {Proof.} \quad}
\newcommand{\fproof}{\hfill $\square$ \bigskip}

\newtheorem{assum}{A}
\renewcommand{\theequation}{\thesection.\arabic{equation}}
\numberwithin{equation}{section}
\def\argmax{\mathop{\mathrm{argmax}}}
\def\RB{\mathbb{R}}
\def\DB{\mathbb{D}}
\def\FB{\mathbb{F}}
\def\FC{\mathcal{F}}
\def\PC{\mathcal{P}}
\def\EC{\mathbb{R^*}}
\def\AC{\mathcal{A}}
\def\UC{\mathcal{U}}
\def\CC{\mathcal C}
\def\RC{\mathcal R}
\def\HC{\mathcal H}
\def\VC{\mathcal V}
\definecolor{LightCyan}{rgb}{0.88,1,1}
\def\cF{{\mathcal {F}}}
\def\MM{{\mathbb{M}}}
\def\RR{{\mathbb{ R}}}
\def\cB{{\mathcal{B}}}
\def\EE{{\mathbb{ E}}}
\def\coz{\mathrm{coz}}
\def\1B{\text{1\!\!I}}
\def\tY{\tilde{Y}}
\def\tN{\tilde{N}}
\def\tl{\tilde{\lambda}}
\def\tH{\tilde{H}}
\def\tp{\tilde{p}}
\def\tr{\tilde{r}}
\def\tq{\tilde{q}}
\def\tz{\tilde{z}}
\def\tk{\tilde{k}}
\def\hX{\hat{X}}
\def\hZ{\hat{Z}}
\def\hA{\hat{A}}
\def\hY{\hat{Y}}
\def\hu{\hat{u}}
\def\hb{\hat{b}}
\def\hg{\hat{g}}
\def\hc{\hat{c}}
\def\hQ{\hat{Q}}
\def\hH{\hat{H}}
\def\hp{\hat{p}}
\def\hth{\hat{\theta}}
\def\hq{\hat{q}}
\def\hr{\hat{r}}
\def\hxi{\hat{\xi}}
\def\hy{\hat{y}}
\def\hk{\hat{k}}
\def\hz{\hat{z}}
%\def\r{\rangle}
\def\l{\langle}
\def\<{\langle}
\def\>{\rangle}
\def\P{\mathbb{P}}
\def\R{\mathbb{R}}
\def\l{\lambda}
\def\S{\mathcal{S}}
\def\E{\mathbb{E}}
\def\N{\mathbb{N}}
\def\L{\mathcal{L}}
\def\F{\mathcal{F}}
\def\H{\mathbb{H}}
\def\Y{\mathcal{Y}}
\def\G{\mathbb{G}}
\def\reali{\mathbb{R}}
\def\realio{\mathbb{R}_{0}}
\def\L{\Lambda}
\def\d{\diamond}
\def\L{\ell}
\def\d{\diamond}



\begin{document}

\title{Multiparameter L\'evy white noise theory and applications}

\author{Olfa Draouil$^{1}$, Rahma Yasmina Moulay Hachemi$^{2}$ \& Bernt \O ksendal$^{3}$
}
\date{26 September 2026}
\maketitle 
\footnotetext[1]{Department of Mathematics, University of Tunis El Manar, Tunisia\newline
Email: olfa.draouil@fst.utm.tn}


\footnotetext[2]{%
Department of Mathematics, University of Oslo, Norway.\newline
Email: ryhachem@math.uio.no}

\footnotetext[3]{%
Department of Mathematics, University of Oslo, Norway. \\
Email: oksendal@math.uio.no.}
\paragraph{MSC [2020]:}
\emph{30B50; 34A08; 35D30; 35D35; 35K05; 35R11; 60H15; 60H40}

\paragraph{Keywords:}
\emph{L\' evy sheets, compensated Poisson random measures, white noise calculus.}

\begin{abstract}
We construct a white noise theory and white noise calculus for the (multi-parameter) L\' evy sheet and its compensated Poisson random measures. The theory is applied to a stochastic partial differential equation subject to L\' evy noise.
  \end{abstract}

\section{Introduction}

The theory of white noise for Brownian motion was initiated by T.Hida \cite{H} and further developed by him and coauthors Kuo, Potthoff and Streit in \cite{HKPS}. Subsequently an extension to (multiparameter)  Brownian sheet was carried out and applied to stochastic partial differential equations by Holden, Ub\o e, \O ksendal and Zhang  in \cite {HOUZ}. 

Later Di Nunno, \O ksendal \& Proske \cite{DOP1} and \O ksendal \& Proske \cite{OP} developed a white noise theory for L\' evy processes and for compensated Poisson random measures, respectively. See the book \cite{DOP} for a comprehensive coverage of this.  

The theory of (multiparameter) L\' evy sheets was developed by Iafrate \& Ricciutu  in \cite{IR} and by Dalang \& Walsh \cite{DW}. See also the earlier paper Adler et al  \cite{AMSW}. 

Subsequently a white noise theory for L\' evy sheets was carried out in Chapter 5 of the book Holden et al \cite{HOUZ}. 

The purpose of this paper is to develop this theory further, in particular we prove a fundamental relation between Wick products and
Skorohod integrals. Moreover, we apply this theory to study a stochastic  equation modeling population growth in a random environment, driven by both Brownian and
L\' evy multiparameter white noise.

The content of our paper is the following:
\begin{itemize}
    \item 
    In Section 2 we recall the theory of L\' evy sheets and we give a construction of L\' evy sheets based on the Bochner-Minlos theorem in the white noise probability space.
    \item 
    In Section 3 we present a white noise theory for multiparameter L\' evy sheets.
    \item 
    In Section 4 we develop a white noise calculus based on the Wick product
    \item 
    Finally, in Section 5 we apply the theory to study a time-fractional stochastic partial differential equation modeling population growth in a Brownian-L\' evy noisy environment.
\end{itemize}


\section{Multiparameter L\' evy sheets and L\' evy White Noise}

%We refer to Section \ref{new} for another example with detailed numerics and figures.

\subsection{Multiparameter L\'{e}vy sheets}
%For the convenience of the reader we include a review of the main ideas and concepts of the %construction in Chapter 5 of \cite{HOUZ}.


We recall the theory of L\' evy sheets, mainly based on \cite{IR}, \cite{DW}:

\begin{definition}
  A \emph{multiparameter L\'{e}vy sheet} is a stochastic process $L(x)=L(x,\omega); x \in \R_+^d,\omega \in \Omega$ defined on a probability space $(\Omega,\F,\P)$ and satisfying the following conditions:
\begin{itemize}
\item
    $L$(0)=0 a.s.
\item
$L$ has stationary, independent increments with respect to the natural partial ordering on $\R_+^d$
\item
$L$ is continuous in probability and has c\'adl\'ag paths a.s. 
\end{itemize}
\end{definition}

\begin{remark}
If $R=\Pi_{j=1}^n (x_j,x'_j]$ is a box in $\R_+^d$ then $\Delta_R L$ denotes the increment of $L$ over the box $R$. For example, if $d=2$ then
$$\Delta_R L= L(x_1',x_2')-L(x_1',x_2) - L(x_1, x_2') + L(x_1,x_2).$$
\end{remark}
\begin{remark}    
If we also assume that $L$ has continuous paths, then $L$ becomes a \emph{Brownian sheet}.
\end{remark}


\begin{definition}
\begin{myenumerate}
    \item
    The jump of $L$ at the point $x \in \R_{+}^d$ is defined by
    \begin{align}
        \Delta L_x= L(x)-L(x^{-}).
    \end{align}
    \item
Let $R$ denote a closed finite box in $\R_+^d$ not containing 0, and let $U$ be a Borel subset of $\R_0:=\R \setminus 0$. Note that because the paths are c\'adl\'ag there are only finitely many jumps in $R$
Define \emph{the jump measure} $N(R,U)$ to be the number of jumps of $L$ within the box $R$ of jumps size in $U$. By stationarity the map $R \mapsto N(R,U);  U \in \mathcal{B}(\R_0)$ (the family of all Borel subsets of $\R_0$) depends only on the volume of $R$ and we can write the random measure generated by $(R,U) \mapsto N(R,U)$ in differential form as  $N(dx, du)$, where $dx$ and $du$ are the Lebesgue measure differentials corresponding to $R$ and $U$, respectively.
\item
Put $R_1=[0,1]^d$. The measure $U \mapsto \nu(U):=\E [N(R_1,U)];  U \in \mathcal{B}(\R_0)$ is called the L\'{e}vy measure of $L$.
\item
The corresponding \emph{compensated jump measure} is defined by
$$ \widetilde{N}(dx,dz)= N(dx,dz) - \nu(dz) dx,$$
where $dx=dx_1 dx_2 ... dx_n$ is the volume differential at $x=(x_1, x_2, ... ,x_d) \in \R^d$.
\end{myenumerate}
\end{definition}
\subsection{The L\'{e}vy white noise probability space and the pure jump L\'{e}vy sheet}

In the following we let $\S=\S(\R^d)$ denote the Schwartz space of rapidly decreasing smooth functions on $\R^d$, and we let $\S'=\S'(\R^d))$ denote its dual, which is the space of tempered distributions on $\R^d$, equipped with the weak$^{*}$ topology. 

It is possible to construct a \emph{pure jump L\'{e}vy sheet} $L$ corresponding to a given  L\'{e}vy measure $\nu$  
as a process on the Borel $\sigma$-algebra $\cF = \cB(\S'(\R^d))$. The details are the following:


Let $\nu$ be a measure on $\R_0$ such that
\begin{align}
\int_{\R_0} z^2 \nu(dz) < \infty.
\end{align}
By the Bochner-Minlos-Sazonov theorem\index{Bochner-Minlos-Sazonov theorem} (see \cite{GV}), there exists a probability measure $\P$ on the Borel $\sigma$-algebra $\cF = \cB(\S'(\R^d))$ such that
\begin{equation}
\int_\Omega e^{i\langle\omega,\varphi \rangle}\P(d\omega)=\exp\left(\int_{\R^n}\Psi(\varphi(x))dx\right);
\quad \varphi\in\mathcal{S}(\R^d), \label{boch}
\end{equation}
where $i=\sqrt{-1}$ and $\langle\omega,\varphi\rangle$ denotes the action of $\omega\in \Omega:=\mathcal{S}'(\R^d)$ on $\varphi\in \mathcal{S}(\R^d)$ and 
\begin{equation}
\Psi(w)=\int_{\R}\left(e^{i\,w\,z}-1-i\,w\,z\right)\nu(dz); w \in \R.
\end{equation}

We let $\E[\cdot]=\E_{\P}[\cdot]$ and $Var[\cdot]=Var_{\P}[\cdot]$ denote expectation and variance, respectively, with respect to $\P$.\\
Note that by \eqref{boch} we have, for all parameters $a \in \R$,
\begin{align}
    \E| e^{i\langle \omega,a \varphi \rangle)}]=\exp\left(\int_{\R^n}\int_{\R_0} \{ e^{i a \varphi(x)z} - 1 - i a \varphi(x)z\} \nu(dz) dx \right) \label{boch2}
\end{align}
Expanding the exponential function on both sides of \eqref{boch2} and comparing the terms with the same power of $a$, we obtain the following result:
\begin{lemma}
\label{2.5}
Let $\varphi\in\mathcal{S}(\R^d)$. 
Then we have
\begin{equation}
\E[\langle\cdot,\varphi\rangle]=0
\end{equation}
and
\begin{equation}
Var[\langle\cdot,\varphi \rangle]:=\E[\langle\cdot,\varphi\rangle^2] = M \int_{\R^d} \varphi^2(x)dx, \label{Var}
\end{equation}
where
\begin{align}
    M=\int_{\R_0} z^2 \nu(dz).
\end{align}
Therefore,
\begin{align}
    \E[(\langle \omega,\varphi_1\rangle - \langle\omega,\varphi_2\rangle)^2]=M \int_{\R^d} (\varphi_1(x)-\varphi_2(x))^2 dx=M \| \varphi_1 - \varphi_2 \|^2_{L^2(\lambda)}, 
\end{align}
where $\lambda$ is the Lebesgue measure on $\R^d$.

\end{lemma}

Using this result, we can for each  $\omega \in \Omega$ extend the definition of $\langle \omega,f\rangle$ from $f\in\mathcal{S}(\R^d)$ to any $f\in L^2(\lambda)$ as follows:

If $f\in L^2(\lambda)$, choose $f_n\in\mathcal{S}(\R^n)$ such that $f_n\rightarrow f$ in $L^2(\lambda)$. 
Then by Lemma \ref{2.5} we see that $\{\langle
\omega,f_n\rangle\}_{n=1}^\infty$ is a Cauchy sequence in $L^2(\P)$ and hence convergent in $L^2(\P)$. 
Moreover, the limit depends only on
$f$ and not the sequence $\{f_n\}_{n=1}^\infty$. 
We denote this limit by $\langle \omega,f\rangle; \  \omega \in \Omega, f \in L^2(\lambda).$

In particular, if we define
$$L(x):=L(x,\omega):=\langle\omega,\chi_
{\Pi_{j=1}^d [0,x_j]} \rangle;\quad x=(x_1, x_2, ... ,x_d) \in \R_+^d,$$
then we have the following result:

\begin{theorem}
The stochastic process $L(x)$, $x \in \R_+^d$ has a c\`{a}dl\`{a}g version, also denoted by $L$. 
This process $L(x)\ ; x\in \R_+^d$ is a (pure jump) L\'{e}vy
sheet with L\'{e}vy measure $\nu$.
\end{theorem}


\section{Multiparameter L\' evy white noise theory}
Our presentation in this section is based on a combination of the results in the papers \cite{IR}, \cite{DW}, \cite{DOP} and \cite{OP} and the book \cite{HOUZ}.\\

In the following we will for simplicity assume that
\begin{align}
\int_{\R_0} z^2 \nu(dz) < \infty.
\end{align}
Then by the L\'{e}vy-Khinchine theorem the L\'{e}vy sheet $L$ can be given the following representation:
\begin{align}
    L(x)=\alpha \Pi_{j=1}^n x_j + \sigma B(x) + \int_{\R_0} z \widetilde{N}(x, dz); \ x \in \R_+^d,
\end{align}
for some constants $\alpha, \sigma$, where $B(x)$ is a Brownian sheet.

We will from now on assume that $\alpha=\sigma=0$. 
Then we have a \emph{pure jump L\'evy sheet}, with representation
\begin{align} 
L(x)=\int_{\R_0} z \widetilde{N}(x,dz)=\int_0^x \int_{\R_0} z \widetilde{N}(d\xi,dz); \ x \in \R_+^d. \label{LN}
\end{align}
In differential form this becomes
\begin{align}\label{ldx}
    L(dx)= \int_{\R_0} z \widetilde{N}(dx,dz).
\end{align}


We proceed to define the white noise of the L\'{e}vy sheet we constructed in the previous section. We will follow the presentation in \cite{DOP} Section 13.2  closely, but with the necessary modifications due to the fact that we are now dealing with an $n$-parameter L\'{e}vy sheet, not a classical 1-parameter L\'{e}vy process. Here are the details:\\

From now on we assume that the L\'{e}vy measure $\nu $ satisfies the
following condition:\\
For all $\varepsilon >0$ there exists $\lambda >0$ such that
\begin{equation}
\int_{\realio\backslash (-\varepsilon ,\varepsilon )}\exp (\lambda
\left\vert z\right\vert )\nu (dz)<\infty   \label{nu}
\end{equation}
One can prove that if this holds, then $\nu $ has finite moments of all orders $n\geq 2$. 
For example, the condition clearly holds if $\nu $ is supported on a compact subset of $\R$.

From this condition one can prove that the polynomials are dense in $L^{2}(\rho )$,
where
\begin{equation}
\rho (dz)=z^{2}\nu (dz).  
\end{equation}
See \cite{NualartSchoutens}. We let $\left\{\eta_{m}\right\} _{m\geq 0}=\left\{ \eta_0 =1,\eta_{1},\eta_{2},...\right\} $ \label{simb-1303} be the orthogonolization of
$\left\{1,z,z^{2},...\right\} $ with respect to the inner product of $L^{2}(\rho )$.

Define
\begin{equation}
p_{j}(z):=\left\Vert \eta_{j-1}\right\Vert _{L^{2}(\rho )}^{-1}z\eta_{j-1}(z);
\text{ }j=1,2,...  \label{10.3}
\end{equation}
and
\begin{equation}
m_{2}:=\left( \int_{\realio} z^{2}\nu (dz)\right) ^{\frac{1}{2}}
=\left\Vert \eta_{0}\right\Vert _{L^{2}(\rho )}
=\left\Vert 1\right\Vert_{L^{2}(\rho )}.  \label{10.4}
\end{equation}
In particular,
\begin{equation}
p_{1}(z)=m_{2}^{-1}z\text{ or }z=m_{2}p_{1}(z).  \label{10.5}
\end{equation}
Then $\left\{ p_{j}(z)\right\} _{j\geq 1}$ is an \emph{orthonormal basis}
for $L^{2}(\rho )$.



Recall that the {\it Hermite polynomials\/} $h_n(x); x \in \R$ are defined by
$$h_n(x)=(-1)^n
e^{{1}/{2}x^2}\frac{d^n}{dx^n}(e^{-{1}/{2}x^2});\quad
n=0,1,2,\cdots.\leqno{(2.2.1)}$$
Thus the first Hermite polynomials are
\begin{align}
h_0(x)&=1,\;h_1(x)=x,\;h_2(x)=x^2-1,\;h_3(x)=x^3-3x,\cr
h_4(x)&=x^4-6x^2+3,\;h_5(x)=x^5-10x^3+15x,\cdots.
\end{align}
The \emph{Hermite functions} $\xi_n(x)$ are defined by
$$\xi_n(x)=\pi^{-{1}/{4}}((n-1)!)^{-{1}/{2}}e^{-{1}/{2}x^2}h_{n-1}
(\sqrt{2}x);\quad n=1,2,\cdots.$$
We note the following  useful properties of $\xi_n$:
\begin{itemize}
\item
$\xi_n\in{\cal S}(\R)\quad\text{for all}\;n.$
\item
 The collection $\{\xi_n\}^\infty_{n=1}$ constitutes an
orthonormal basis for $L^2(\R).$
\item
$\sup\limits_{x\in\R}|\xi_n(x)|=\mathcal{O}(n^{-{1}/{12}}).$
\end{itemize}
See Hille and
Phillips (1957), Chapter 21, for proofs of these statements.
\vskip 0.2cm
Using these functions we define an orthogonal basis for $L^2(\P)$ as follows:\\
Let
$\beta=(\beta_1,\cdots,\beta_d)$ denote $d$-dimensional multi-indices with
$\beta_1,\cdots,\beta_d\in\N$, where $\mathbb{N}$ denotes the set of all natural numbers. By the above it follows that the family of tensor
products
$$\xi_\beta:=\xi_{(\beta_1,\cdots,\beta_d)}:=\xi_{\beta_1}\otimes\cdots
\otimes
\xi_{\beta_d}\;;\;\beta\in\N^d$$
forms an orthonormal basis for $L^2(\R^d)$. Let
$\beta^{(j)}=(\beta_1^{(j)},\beta_2^{(j)},\cdots,\beta_d^{(j)})$ be
the $j$th multi-index in some fixed ordering of all $d$-dimensional
multi-indices $\beta=(\beta_1,\cdots,\beta_d)\in\N^d$. We 
may assume that this ordering has the property that
$$i<j\Rightarrow\beta_1^{(i)}+\beta_2^{(i)}+\cdots+\beta_d^{(i)}\leq
\beta_1^{(j)}+\beta_2^{(j)}+\cdots+\beta_d^{(j)},$$
i.e., that the $\{\beta^{(j)}\}^\infty_{j=1}$ occur in increasing order.


Next, define the bijection $\kappa :\mathbb{N}\times \mathbb{N}\longrightarrow
\mathbb{N}$\label{simb-new2} by
\begin{equation}
\kappa (i,j)=j+(i+j-2)(i+j-1)/2;\quad  i,j=1,2, ...
\end{equation}
%\[
%\begin{array}{cccccccccc}
%(1) &  & (2) &  & (4) &  & (i) &  &  &  \\
%\bullet  & \longrightarrow  & \bullet  &  & \bullet  & \cdots  & \bullet  &
%\longrightarrow  &  &  \\
%(3) & \swarrow  & (5) & \swarrow  &  &  &  &  &  &  \\ 
%\bullet  &  & \bullet  &  &  &  &  &  &  &  \\ 
%(6) & \swarrow  &  &  &  &  &  &  &  &  \\ 
%\bullet  &  &  &  &  &  &  &  &  &  \\ 
%\vdots  &  &  &  &  &  &  &  &  &  \\
%(j) &  &  &  &  &  &  &  &  &  \\ 
%\bullet  &  &  &  &  &  &  &  &  &  \\ 
%\downarrow  &  &  &  &  &  &  &  &  & 
%\end{array}
%\]
For $x=(x_1, x_2, ... , x_d) \in \R^d$ we proceed as in in Section 2.2.1 in \cite{HOUZ} and define $\left\{ e _{j}(x)\right\} _{j\geq 1}$ to be a fixed ordering of all $d$-times tensor products of the Hermite functions.
This family forms an orthonormal basis of $L^2(\R^d).$
Put
\begin{equation}
\theta _{\kappa (j,k)}(x,z):=e _{j}(x)p_{k}(z); \  x\in \R^n, z \in \R_0.
\end{equation}

In the following we let $\mathcal{J}$ denote the family of all finite multi-indices $\alpha=(\alpha_1, \alpha_2, ..., \alpha_m); \ m=1,2, ...$ of nonnegative integers $\alpha_j$. 

If $\alpha \in \mathcal{J}$ with $Index(\alpha )=j$  (the index of the last non-zero element) and $\left\vert \alpha
\right\vert =m$, we define $\theta ^{\otimes \alpha }: (\R^d \times \R_0)^m \mapsto \R$ by
\begin{eqnarray}
&&\theta ^{\otimes \alpha }(x^{(1)},z_{1};...;x^{(m)},z_{m}) \\
&=&\theta _{1}^{\otimes \alpha _{1}}\otimes ...\otimes \theta _{j}^{\otimes
\alpha _{j}}(x^{(1)},z_{1};... ;x^{(m)},z_{m})  \nonumber \\
&=&\underset{\alpha _{1}\text{ factors}}{\underbrace{\theta
_{1}(x^{(1)},z_{1})\cdot ...\cdot \theta _{1}(x^{(\alpha _{1})},z_{\alpha _{1}})}
}\cdot ...\cdot \underset{\alpha _{j}\text{ factors}}{\underbrace{\theta
_{j}(x^{(m-\alpha _{j}+1)},z_{m-\alpha _{j}+1})\cdot ...\cdot \theta
_{j}(x^{(m)},z_{m})}}.  \nonumber
\end{eqnarray}
We put $\theta _{i}^{\otimes 0}=1.$\\
Finally we let $\theta ^{\hat{\otimes }\alpha }$ denote the
\emph{symmetrized} tensor product\index{tensor product!symmetrized} of the $\theta _{k}$ $^{\prime }s:$
\begin{equation}
\theta ^{\hat{\otimes }\alpha }(x^{(1)},z_{1};...;x^{(m)},z_{m})
=\theta_{1}^{\hat{\otimes }\alpha _{1}}\otimes ...\otimes
\theta _{j}^{\hat{\otimes }\alpha _{j}}(x^{(1)},z_{1};...;x^{(m)},z_{m}).  
\end{equation}
The symmetrization is with respect to the $m$ pairs $(x^{(1)},z_1), (x^{(2)}z_2), \cdots (x^{(m)}z_m)$.
We set $K_0=1$ and for general $\alpha \in \mathcal{J}$ we define
\begin{equation}
K_{\alpha }:= I_{\left\vert \alpha \right\vert }\left(
\theta ^{\hat{\otimes }\alpha }\right),
\end{equation}
where in general $I_m$ is the symmetrized iterated integral with respect to $\widetilde{N}$. This integral is defined for symmetric functions $g: (\R^d \times \R_0)^m \mapsto \R$ by
\begin{align}
   I_m(g) &= \int_{(\R^d \times \R_0)^m } g(x^{(1)},z_1; x^{(2)},z_2; ... ;x^{(m)},z_m) \widetilde{N}^{\otimes m}(dx^{(1)} dz_1; ..., dx^{(m)} dz_m)\nonumber\\
   &:= (m!)^{d} J_m(g), 
\end{align}
where $J_m(g)$ is the $m$ times iterated integral
\begin{align}
    &J_m(g)\nonumber\\
    &=\int_{\R^d}\int_{\R_0} \Big(\int_{-\infty}^{x^{(m)}} \int_{\R_0}\Big( ... \Big(\int_{-\infty}^{x^{(2)}}g(x^{(1)},z_1; x^{(2)},z_2\cdots \nonumber\\
    &\cdots x^{(m)},z_m) ) \widetilde{N}(dx^{(1)},dz_1)\Big)  ... \Big) \widetilde{N}(dx^{(m)},dz_m). 
\end{align}
Here we are using the following notation, for $x=(x_1,x_2, ... ,x_m) \in \R^m$,
\begin{align}
    \int_{-\infty}^x= \int_{-\infty}^{x_m} \int_{-\infty}^{x_{m-1}}  ... \int_{-\infty}^{x_1}. 
\end{align}


\begin{example}
    For example, if $m=1$ we get
    \begin{align}
        I_1(g)=\int_{\R^d} \int_{\R_0} g(x,z) \widetilde{N}(dx,dz),
    \end{align}
    and if $m=2$ we have
    \begin{align}
        I_2(g)&=\int_{\R^d} \int_{\R_0} \Big( \int_{\R^d} \int_{\R_0} g(x^{(1)},z_1; x^{(2)},z_2) \widetilde{N}(dx^{(2)},dz_2) \Big)\widetilde{N}(dx^{(1)},dz_1)\nonumber\\
        &=2^d\int_{\R^d}\int_{\R_0} \Big(\int_{-\infty}^{x^{(2)}} \int_{\R_0} g(x^{(1)},z_1; x^{(2)},z_2) \widetilde{N}(dx^{(1)},dz_1)\Big) \widetilde{N}(dx^{(2)},dz_2)
    \end{align}
\end{example}
\begin{example}
\label{Ex10.1}
With $\varepsilon ^{(k)}=(0,...,0,1,0,...)$ with $1$ on $k$ th
place, we have, writing $\varepsilon ^{(\kappa (i,j))}=\varepsilon ^{(i,j)}$,
\begin{equation}
K_{\varepsilon ^{(i,j)}}
= I_1(\theta^{\otimes \varepsilon^{(i,j)}})
=I_{1}\left( \theta_{\kappa (i,j)} \right)
=I_{1}\left( e _{i}(x)p_{j}(z)\right) .  \label{10.11}
\end{equation}%

As in the Brownian sheet case one can prove that $\{ K_\alpha  \}_{\alpha \in \mathcal{J}}$ are
orthogonal in $L^2(\P)$ and
\[
\Vert K_\alpha \Vert^2_{L^2(\P)} = \alpha!.
\]
Note that if $|\alpha|=m$,
\[
\alpha ! =\Vert K_\alpha \Vert^2_{L^2(\P)} = m! \Vert \theta^{\hat\otimes \alpha} \Vert^2_{L^2((\lambda\times\nu)^m)},
\]
where as before $\lambda$ denotes Lebesgue measure on $\R^d$.
By our construction of $\theta ^{\hat{\otimes }\alpha }$ we know that
any $f\in \widetilde{L}^{2}((\lambda \times \nu )^{m})$ can be written%
\begin{equation}
f(x^{(1)},z_{1};...; x^{(m)},z_{m})=\sum_{\left\vert \alpha \right\vert
=m}c_{\alpha }\theta ^{\hat{\otimes }\alpha
}(x^{(1)},z_{1},..., x^{(m)},z_{m}).
\end{equation}
Hence
\begin{equation}
I_{m}(f_{m})=\sum_{\left\vert \alpha \right\vert =m}c_{\alpha }K_{\alpha }
\end{equation}
\end{example}
This gives the following chaos expansion:
\begin{theorem}
\label{Th10.2}
{\bf (Chaos expansion)}
\newline Any $F\in L^{2}(\P)$
has a unique expansion of the form
\begin{equation}
F=\sum_{\alpha \in \mathcal{J}}c_{\alpha }K_{\alpha }= \sum_{m=0}^\infty I_m(f_m).
\label{10.14}
\end{equation}
with
\begin{equation}
f_m=f_m (x^{(1)},z_{1};\cdots ; x^{(m)},z_{m})=\sum_{\left\vert \alpha \right\vert
=m}c_{\alpha }\theta ^{\hat{\otimes }\alpha
}(x^{(1)},z_{1},\cdots ; x^{(m)},z_{m}).
\end{equation}
and $c_{\alpha }\in \mathbb{R}$. 
Moreover,
\begin{equation}
\left\Vert F\right\Vert _{L^{2}(\P)}^{2}=\sum_{\alpha \in \mathcal{J}}\alpha
!c_{\alpha }^{2}.  \label{10.15}
\end{equation}
\end{theorem}

\begin{example}
\label{Ex10.3}
Let $h\in L^{2}(\mathbb{R})$ and define, for $x \in \R_+^d$,
\begin{align}
F=\int_0^{x} h(\xi)L(d\xi).
\end{align}
By  \eqref{LN} this can be written
\begin{align}
    F=I_{1}(h(\xi)\chi_{\Pi_{i=1}^d [0,x_i]}(\xi)z).
\end{align}
\end{example}
Therefore $F$ has the following expansion, writing $\chi_{\Pi_{i=1}^d [0,x_i]}=\chi_{[0,x]}$ for short:
\begin{eqnarray}
F &=&I_{1}(\sum_{i\geq 1}(h\chi_{[0,x]},e _{i})e _{i}(\xi)z)=\sum_{i\geq 1}(h\chi_{[0,x]},e
_{i})I_{1}(e _{i}(\xi)z) \\
&=&\sum_{i\geq 1}\left( \int_0^x h(\xi)e _{i}(\xi)d\xi\right)
K_{\varepsilon ^{(i,1)}}m_{2}.  \nonumber
\end{eqnarray}
Recall that $z=m_{2}p_{1}(z);$ see (\ref{10.5}).
In particular, choosing $h=1$ we get the following expansion for the L\' evy sheet $L$:
\begin{equation}
L(x)=\sum_{i\geq 1}\left( \int_{0}^{x}e _{i}(\xi)d\xi\right)
K_{\varepsilon ^{(i,1)}}m_{2}.  \label{LK}
\end{equation}

\subsection{The Skorohod integral}
The following definition of the {\it Skorohod integral} with respect to the L\'{e}vy field $L(x)$ is an extension to the  multiprameter case of a definition originally due to Kabanov
(1975):
\bigskip

\begin{definition}[The Skorohod integral]
\label{Skoro}
Let $Y(x)$, $x\in\mathbb{R}^d$, be a measurable stochastic process
such that
\[
\mathbb{E}[Y^2(x)]<\infty,
\qquad x\in\mathbb{R}^d.
\]
Then, for each $x\in\mathbb{R}^d$, $Y(x)$ admits a chaos expansion
of the form
\[
Y(x)=\sum_{n=0}^{\infty}I_n(f_n(\cdot,x)),
\]
where
\[
f_n(\cdot,x)\in
\widehat{L}^2\big((\lambda\times\nu)^n\big),
\qquad n\geq 1,
\]
and
\[
I_0(f_0(\cdot,x))=\mathbb{E}[Y(x)].
\]

For $n\geq 0$, let $\widetilde{f}_n$ denote the symmetrization, with
respect to all $n+1$ variables, of the function
\[
\begin{aligned}
&(x^{(1)},z_1,\ldots,x^{(n+1)},z_{n+1})
\\
&\qquad\longmapsto
z_{n+1}
f_n(x^{(1)},z_1,\ldots,x^{(n)},z_n,x^{(n+1)}).
\end{aligned}
\]

Suppose that
\[
\sum_{n=0}^{\infty}
(n+1)!
\|\widetilde{f}_n\|^2_{
L^2((\lambda\times\nu)^{n+1})}
<\infty.
\]
Then $Y(\cdot)$ is said to be \emph{Skorohod integrable} with respect
to $L$, and its Skorohod integral is defined by
\begin{equation}
\label{Skor}
\int_{\mathbb{R}^d}Y(x)L(dx)
:=
\sum_{n=0}^{\infty}
I_{n+1}(\widetilde{f}_n).
\end{equation}
\end{definition}

Note that 

$${\rm
\E}\left[\left(\int_{\R^d}Y(x)L(dx)\right)^2\right]=\sum_{n=0}^\infty(n+1)!||\tilde{f}_n||^2_{L^2((\lambda^d\times\nu)^{n+1})}<\infty$$
and therefore $\int_{\R^d}Y(x)L(dx)\in L^2(\P)$. Moreover, 

$$\E \left[\int_{\R^d}Y(x)L(dx)\right]=0.$$




\subsection{The L\'{e}vy-Hida spaces and the Kondratiev spaces}
We now introduce some useful stochastic test function spaces and stochastic distribution spaces.\\

In the following we put
\begin{equation}
(2\mathbb{N})^{k\alpha }=\prod_{j\geq 1}(2j)^{k\alpha _{j}},
\end{equation}%
if $\alpha =(\alpha _{1,}\alpha _{2,},...)\in \mathcal{J}$.






\begin{definition}
 (The L\'{e}vy-Hida-Kondratiev spaces of stochastic test functions and stochastic
distributions)

\parindent5mm
\item{{\bf a)}} {\it (The Kondratiev stochastic test function spaces)}. For $0\le\rho\le 1, k=1,2,...$, let $({\cal S})_{\rho,k}=({\cal
S})_{\rho,k}^{(L)}$ consist of those
$$\phi=\sum_{\alpha\in{\cal J}} c_\alpha K_\alpha(\omega)\in L^2(\P)\hskip1cm (c_\alpha\in\R\ \hbox{constants})$$
such that

$$||\phi||^2_{\rho,k}:=\sum_{\alpha\in{\cal J}}c_\alpha^2(\alpha!)^{1+\rho}(2\N)^{k\alpha}<\infty\hskip1cm \hbox{for {\it all} 
}k\in\N $$

\item{{\bf b)}} {\it (The Kondratiev stochastic distribution spaces)}. For $0\le\rho\le 1, k=1.2...$ let $({\cal S})_{-\rho,-k}=({\cal
S})_{-\rho,-k}^{(L)}$ consist of all formal expansions
$$F=\sum_{\alpha\in{\cal J}}b_\alpha K_\alpha(\omega)$$
such that

$$||F||^2_{-\rho,-k}:=\sum_{\alpha\in{\cal J}}b_\alpha^2(\alpha!)^{1-\rho}(2\N)^{-k\alpha}<\infty\hskip1cm \hbox{for {\it some}
}k\in\N$$


We define 
$$({\cal S})_\rho=\cap_{k=1}^{\infty} \S_{\rho,k},$$
equipped with the projective topology (intersection) defined by the norms $||\cdot||_{\rho,k};\
k=1,2,\dots$ and we define 
$$(\S)_{-\rho}=\cup_{k=1}^{\infty}(\S)_{-\rho,-k},$$ 
equipped with the inductive topology (union) defined by the norms 
$||\cdot||_{-\rho,-k};\ k=1,2,\dots$. Then $({\cal S})_{-\rho}$ becomes the dual of $({\cal S})_\rho$ and the action of $F=\sum
b_\alpha K_\alpha\in ({\cal S})_{-\rho}$ on $\phi=\sum a_\alpha K_\alpha\in ({\cal S})_{\rho}$ is given by

$$<F,\phi>=\sum_\alpha a_\alpha b_\alpha \alpha!\leqno{(5.3.30)}$$
\item{{\bf c)}} The space $({\cal S})=({\cal S})_0^{(L)}$ is called the {\it L\' evy-Hida stochastic test function space} and the space
$({\cal S})^*=({\cal S})_{-0}^{(L)}$ is called the {\it L\' evy-Hida distribution space}.\\

{\bf Generalized expectation.}\newline
If  $F= \sum_{\alpha \in \mathcal{J}} a_{\alpha }K_{\alpha } \in (\mathcal{S})^{* }$, we define
the generalized expectation $\E[F]$ of $F$ by
$$
\E[F] = a_0.
$$ 
Note that $\E[K_{\alpha }] = 0$ for all $\alpha \ne 0$.
Therefore the generalized expectation coincides with the usual expectation if $F \in L^2(\P)$. 

    
\end{definition}

\subsection{White noise for L\' evy sheets and their compensated Poisson random measures}
Now we have the prerequisites to define the white noise of a L\'evy sheet and its associated Poisson random measure. In view of \eqref{LK} the following definition is natural:

\begin{definition} 
\begin{enumerate}
\item[(i)]
The \emph{L\'{e}vy white noise}
$$\overset{\bullet }{L
}(x):=V(x):= \frac{\partial^d L(x)}{\partial x_1 \partial x_2  ... ,\partial x_d},\quad x\in \R^d$$
is defined by the expansion
\begin{eqnarray}
V(x):&=&m_{2}\sum_{i\geq 1}e _{i}(x)K_{\varepsilon
^{(i,1)}}=m_{2}\sum_{i\geq 1}e _{i}(x)I_{1}(e _{i}(x)p_{1}(z))\nonumber\\
&=&m_{2}\sum_{i\geq 1}e _{i}(x)I_{1}(e _{i}(x)z).  
\end{eqnarray}%

\item[(ii)]
The \emph{white noise} $\overset{\bullet }{\widetilde{N}%
}(x,z)$ of the compensated random measure  $\widetilde{N}(dx,dz)$ is defined by the expansion%
\begin{equation}
\overset{\bullet }{\widetilde{N}}(x,z)=\sum_{i,j\geq 1}e
_{i}(x)p_{j}(z)K_{_{\varepsilon ^{(i,j)}}}(\omega ).  
\end{equation}
\end{enumerate}
\end{definition}




\begin{remark}
\label{Rem10.6}Note that in the case of the expansion for $\overset{\bullet }{L}(x); x \in \R^d$ we have%
\[
\sum_{\alpha \in \mathcal{J}}c_{\alpha }^{2}\alpha !(2\mathbb{N})^{-q\alpha
}=m_{2}^{2}\sum_{i\geq 1}e _{i}^{2}(x) 2^{-q}\kappa (i,1)^{-q}<\infty \text{ }
\]% 

for $q\geq 2.$ Here we have used that $\kappa (i,1)=1+(i-1)i/2\geq i,$  and the following
well-known estimate for the Hermite functions:%
\begin{equation}
\sup_{x\in \mathbb{R}}\left\vert e _{m}(x)\right\vert =\mathcal{O}(m^{-\frac{1}{12}%
}).
\end{equation}%
This estimate is well-known for the 1-parameter Hermite functions with $x \in \R$, and it is easy to see that it extends to the tensor products $e_m(x)$ with $x \in \R^n$.
Therefore $\overset{\bullet }{L}(x)\in (\mathcal{S})^{* }$ for all $x$.
 Similarly $\overset{\bullet }{\widetilde{N}}(x,z)\in (\mathcal{S})^{* }$
for all $x,z.$
\end{remark}


\begin{remark}
\begin{enumerate}
\item[(i)]
By comparing the expansions%
\[
L(x)=m_2 \sum_{i\geq 1}\left( \int_{0}^{x_1}\int_0^{x_2} ... \int_0^{x_n} e _{i}(\xi) d\xi\right)
K_{\varepsilon ^{(i,1)}}
\]%
and%
\[
\overset{\bullet }{L}(x)=m_{2}\sum_{i\geq 1}e _{i}(x)K_{\varepsilon
^{(i,1)}},
\]%
we get formally
\begin{equation}
\overset{\bullet }{L}(x)=\frac{\partial^n L(x)}{\partial x_1 \partial x_2 ..., \partial x_n}\text{ (derivative in }(%
\mathcal{S})^{* }). 
\end{equation}
This can be proved rigorously.
%See Problem \ref{newexe 10.1}.
\item[(ii)] 
Choose a Borel set $\mho $ such that $\overline{\mho }\subset \mathbb{%
R}\backslash \{0\}.$ Then%
\begin{eqnarray*}
\widetilde{N}(x,\mho ) &=&I_{1}\left( \chi _{[0,x]}(\xi)\chi _{\mho
}(z)\right)  \\
&=&\sum_{i,j\geq 1}(\chi _{[0,x]},e _{i})_{L^{2}(\mathbb{R}^n)}(\chi
_{\mho },p_{j})_{_{L^{2}(\nu )}}I_{1}\left( e _{i}(\xi)p_{j}(\zeta)\right) \\
&=&\sum_{i,j\geq 1}\int_{0}^{x}e _{i}(\xi)d\xi\cdot \int_{\mho }p_{j}(z)\nu
(z)K_{\varepsilon ^{(i,j)}}(\omega ).
\end{eqnarray*}% 
This justifies the identity
\begin{equation}
\overset{\bullet }{\widetilde{N}}(x,z)=\frac{\widetilde{N}(dx,dz)}{dx\times
\nu (dz)}\text{ (Radon-Nikodym derivative).}  \label{10.26}
\end{equation}%
\item[(iii)]
Also note that $\overset{\bullet }{L}$ is related to $%
\overset{\bullet }{\widetilde{N}}$ by%
\begin{equation}
\overset{\bullet }{L}(x)=\int_{\realio} \overset{\bullet }{\widetilde{N%
}}(x,z)z\nu (dz).  
\end{equation}%
\dproof
To see this we consider%
\begin{eqnarray*}
&&\int_{\realio}\overset{\bullet }{\widetilde{N}}(x,z)z\nu (dz) \\
&=&\int_{\realio}\sum_{i,j\geq 1}e _{i}(x)p_{j}(z)K_{_{\varepsilon
^{(i,j)}}}(\omega )z\nu (dz)
=\sum_{i\geq 1}e _{i}(x)\sum_{j\geq 1}I_{1}( e
_{i}p_j)\int_{\realio} p_{j}(z)z\nu (dz))  \\
&=&m_2\sum_{i\geq 1}e _{i}(x)I_{1}\left( e _{i}(\xi)z\right) =\overset{%
\bullet }{L}(x).
\end{eqnarray*}
\fproof
\end{enumerate}
\end{remark}


\section{White noise calculus}
\subsection{The Wick product}

The Wick product is important in the study of
stochastic (ordinary and partial) differential equations. In
general, one can say that the use of this product corresponds to
-- and extends naturally -- the use of It\^{o} integrals. We now
explain this in more detail.

\begin{definition}The {\it Wick product} $F\diamond G$ of two elements
$$F=\sum\limits_\alpha a_\alpha K_\alpha,\;G=\sum\limits_{\beta} b_{\beta}
K_{\beta}\in({\cal S})^*\quad\text{with}\quad
a_{\alpha},b_{\beta}\in\R$$
is defined by
$$F\d 
G=\sum\limits_{\alpha,\beta}a_{\alpha}b_{\beta}K_{\alpha+\beta}.$$
\end{definition}
\begin{example}
 Choose 
 \begin{align}
 &F=L(x)=I_{1}(\sum_{i\geq 1}(\chi_{[0,x]},e _{i})e _{i}(\xi)z)\\
&=\sum_{i\geq 1}(\chi_{[0,x]},e
_{i})I_{1}(e _{i}(\xi)z)=\sum_{i\geq 1}\left( \int_0^x e _{i}(\xi)d\xi\right)
K_{\varepsilon ^{(i,1)}}m_{2}\\
&\text{ and }\\
&G=L(y)=I_{1}(\sum_{i\geq 1}(\chi_{[0,y]},e _{i})e _{i}(\xi)z)\\
&=\sum_{i\geq 1}(\chi_{[0,y]},e
_{i})I_{1}(e _{i}(\xi)z)=\sum_{i\geq 1}\left( \int_0^y e _{i}(\xi)d\xi\right)
K_{\varepsilon ^{(i,1)}}m_{2}.  
\end{align}
Then
\begin{align*}
    &L(x) \diamond L(y)= \sum_{i,j} ( \int_0^x e _{i}(\xi)d\xi)( \int_0^y e _{j}(\xi)d\xi) K_{\epsilon^{(i,1)}+\epsilon ^{(j,1)}}\\
    &=\sum_{i,j} ( \int_0^x e _{i}(\xi)d\xi)( \int_0^y e _{j}(\xi)d\xi) K_{\epsilon^{(i.1)+(j,1)}}\\
    &=\sum_{i,j} ( \int_0^x e _{i}(\xi)d\xi)( \int_0^y e _{j}(\xi)d\xi)I_2((e_ip_1)\widehat{\otimes} (e_j p_1)) \\
    &=\sum_{i,j} ( \int_0^x e _{i}(\xi)d\xi)( \int_0^y e _{j}(\xi)d\xi)\int_{\R}\int_{\R_0} \Big(\int_{\R} \int_{\R_0} e_i(x^{(1)})z_1e_j(x^{(2)})z_2 \widetilde{N}(dx^{(1)},dz_1)\Big) \widetilde{N}(dx^{(2)},dz_2)
\end{align*}

\end{example}

\subsection{Wick products and Skorohod integration}

We now prove a fundamental relation between Wick products and
Skorohod integrals:

\begin{definition}
\label{definition3.13-added} A function $Y:\mathbb{R}^d \to 
(\mathcal{S})^{*}$ (also called an $(\mathcal{S})^{*}$-valued process) is
\emph{$(\mathcal{S})^{*}$-integrable} if
\[
< Y(\cdot),f> \in L^{1}(\mathbb{R}^d), \quad\text{ for all}\quad f\in(\mathcal{S}).
\]
Then the \emph{$(\mathcal{S})^{*}$-integral of $Y$}%
\index{$(\mathcal{S})^{*}$-integral}%
, denoted by $\int_{\mathbb{R}^d} Y(x)dx$, is the (unique) element in
$(\mathcal{S})^{*}$ such that
\begin{equation}
\label{3.68-added}<\int_{\mathbb{R}^d}Y(x)dx, f> = \int_{\mathbb{R}^d} <Y(x),f>
dx, \quad\text{ for all } f\in(\mathcal{S}).
\end{equation}

\end{definition}

\begin{remark}
\label{remark3.14-added} The fact that \eqref{3.68-added} does indeed define
$\int_{\mathbb{R}^d }Y(x)dx$ as an element of $(\mathcal{S})^{*}$ is a
consequence of \cite[Proposition 8.1]{HKPS}.
\end{remark}

Note that if $Y$ is $(\mathcal{S})^{*}$-integrable, then so is $Y\chi_{(a,b]}%
$, for all $a,b\in\mathbb{R}^d$, and we put
\[
\int_{a}^{b} Y(x)dx := \int_{\mathbb{R}^d}Y(x)\chi_{(a,b]}(x)dx.
\]

The following result was first proved by Lindstr\o m et al \cite{LOU} (1992) and subsequently with a different proof by Benth \cite{B} (1993), both in the classical 1-parameter Brownian motion case. Our proof here is an adaptation of Benth's proof in \cite{B} to the multiparameter L\' evy sheet case.

\begin{theorem}[{\cite{LOU}, \cite{B}}]
\label{wick}
Let $Y:\mathbb{R}^d\times\Omega\to\mathbb{R}$ be a Skorohod integrable
stochastic process in the sense of Definition~\ref{Skoro}. Then the
$(\mathcal{S})^{*}$-valued process
\[
x\longmapsto Y(x)\diamond\overset{\bullet}{L}(x)
\]
is $(\mathcal{S})^{*}$-integrable and
\begin{equation}
\label{3.66}
\int_{\mathbb{R}^d}
Y(x)\diamond\overset{\bullet}{L}(x)\,dx
=
\int_{\mathbb{R}^d}Y(x)L(dx),
\end{equation}
where the integral on the right-hand side is the Skorohod integral
defined by
\[
\int_{\mathbb{R}^d}Y(x)L(dx)
:=
\sum_{n=0}^{\infty}
I_{n+1}\big(\widetilde{f}_n\big),
\]
with $\widetilde{f}_n$ denoting the symmetrization of
\[
z_{n+1}f_n
\]
in the $n+1$ variables.
\end{theorem}
\dproof
Let
\[
Y(x)=\sum_{\alpha\in\mathcal J}c_\alpha(x)K_\alpha
     =\sum_{n=0}^{\infty}I_n(f_n(\cdot,x)),
\]
where, for $|\alpha|=n$,
\begin{equation}
f_n(\cdot,x)
=
\sum_{|\alpha|=n}
c_\alpha(x)\,
\theta^{\hat\otimes\alpha}.
\label{fn-expansion}
\end{equation}
By assumption, $Y$ is Skorohod integrable, so that
\[
\sum_{n=0}^{\infty}
(n+1)!
\|\widetilde f_n\|^2_{L^2((\lambda\times\nu)^{n+1})}
<\infty,
\]
where $\widetilde f_n$ denotes the symmetrization of
\[
z_{n+1}f_n
(x^{(1)},z_1,\ldots,x^{(n)},z_n,x^{(n+1)}).
\]

We first compute the Wick integral. Recall that
\[
\overset{\bullet}L(x)
=
m_2\sum_{k=1}^{\infty}
e_k(x)K_{\varepsilon^{(k,1)}},
\]
where
\[
m_2=\left(\int_{\mathbb R_0}z^2\nu(dz)\right)^{1/2}.
\]
Using the definition of the Wick product,
\[
K_\alpha\diamond K_\beta=K_{\alpha+\beta},
\]
we obtain
\begin{align}
Y(x)\diamond\overset{\bullet} L(x)
&=
m_2
\sum_{\alpha\in\mathcal J}
\sum_{k=1}^{\infty}
c_\alpha(x)e_k(x)
K_{\alpha+\varepsilon^{(k,1)}}.
\end{align}
Consequently, by the definition of the $(\mathcal S)^*$-integral,
\begin{align}
\int_{\mathbb R^d}
Y(x)\diamond\overset{\bullet} L(x)\,dx
&=
m_2
\sum_{\alpha\in\mathcal J}
\sum_{k=1}^{\infty}
\left(
\int_{\mathbb R^d}
c_\alpha(x)e_k(x)\,dx
\right)
K_{\alpha+\varepsilon^{(k,1)}} \\
&=
m_2
\sum_{\alpha\in\mathcal J}
\sum_{k=1}^{\infty}
(c_\alpha,e_k)_{L^2(\mathbb R^d)}
K_{\alpha+\varepsilon^{(k,1)}}.
\label{wick-expansion}
\end{align}

We now compute the Skorohod integral using Definition~\ref{Skoro}.
For $|\alpha|=n$, we have
\[
f_n(\cdot,x)
=
\sum_{|\alpha|=n}
c_\alpha(x)\theta^{\hat\otimes\alpha}.
\]
Hence
\[
z_{n+1}f_n
=
\sum_{|\alpha|=n}
c_\alpha(x)
z_{n+1}
\theta^{\hat\otimes\alpha}.
\]
Since
\[
z=m_2p_1(z),
\]
and
\[
\theta_{\kappa(k,1)}(x,z)
=
e_k(x)p_1(z),
\]
we may expand
\[
c_\alpha(x)
=
\sum_{k=1}^{\infty}
(c_\alpha,e_k)_{L^2(\mathbb R^d)}e_k(x)
\]
and therefore obtain
\begin{align}
z_{n+1}f_n
&=
m_2
\sum_{|\alpha|=n}
\sum_{k=1}^{\infty}
(c_\alpha,e_k)_{L^2(\mathbb R^d)}
\,
\theta^{\hat\otimes\alpha}
\otimes
\theta_{\kappa(k,1)}.
\end{align}
After symmetrization,
\begin{align}
\widetilde f_n
&=
m_2
\sum_{|\alpha|=n}
\sum_{k=1}^{\infty}
(c_\alpha,e_k)_{L^2(\mathbb R^d)}
\,
\theta^{\hat\otimes\alpha}
\hat\otimes
\theta_{\kappa(k,1)}.
\end{align}
By the definition of the multi-index notation,
\[
\theta^{\hat\otimes\alpha}
\hat\otimes
\theta_{\kappa(k,1)}
=
\theta^{\hat\otimes
(\alpha+\varepsilon^{(k,1)})}.
\]
Thus
\begin{equation}
\widetilde f_n
=
m_2
\sum_{|\alpha|=n}
\sum_{k=1}^{\infty}
(c_\alpha,e_k)_{L^2(\mathbb R^d)}
\theta^{\hat\otimes
(\alpha+\varepsilon^{(k,1)})}.
\label{tilde-fn}
\end{equation}

By Definition~\ref{Skoro},
\begin{align}
\int_{\mathbb R^d}Y(x)L(dx)
&=
\sum_{n=0}^{\infty}
I_{n+1}(\widetilde f_n)\\
&=
m_2
\sum_{n=0}^{\infty}
\sum_{|\alpha|=n}
\sum_{k=1}^{\infty}
(c_\alpha,e_k)_{L^2(\mathbb R^d)}
I_{n+1}
\left(
\theta^{\hat\otimes
(\alpha+\varepsilon^{(k,1)})}
\right).
\end{align}
Since
\[
K_\beta
=
I_{|\beta|}
\left(
\theta^{\hat\otimes\beta}
\right),
\]
we have
\[
I_{n+1}
\left(
\theta^{\hat\otimes
(\alpha+\varepsilon^{(k,1)})}
\right)
=
K_{\alpha+\varepsilon^{(k,1)}}.
\]
Therefore
\begin{equation}
\int_{\mathbb R^d}Y(x)L(dx)
=
m_2
\sum_{\alpha\in\mathcal J}
\sum_{k=1}^{\infty}
(c_\alpha,e_k)_{L^2(\mathbb R^d)}
K_{\alpha+\varepsilon^{(k,1)}}.
\label{skor-expansion}
\end{equation}

Comparing \eqref{wick-expansion} and \eqref{skor-expansion}, we conclude that
\[
\int_{\mathbb R^d}
Y(x)\diamond\overset{\bullet} L(x)\,dx
=
\int_{\mathbb R^d}Y(x)L(dx).
\]
This proves the result.
\fproof
\begin{remark}
The identity in Theorem~\ref{wick} should be understood with
the Skorohod integral defined by the chaos expansion above. In the
one-parameter case, this integral coincides with the corresponding
It\^o integral for adapted integrands.
In the multiparameter setting, one has to distinguish carefully
between the chaos-defined Skorohod integral and the various
constructions of multiparameter It\^o integrals. The identity
proved above concerns the Skorohod integral defined by the chaos
expansion. The relation between this integral and the
multiparameter It\^o integral requires additional considerations.
The computation on page 229 of
\cite{HOUZ} concerns the latter integral and therefore does not
directly contradict Theorem~\ref{wick}.
\end{remark}
\begin{example}[The one-parameter case]
Let $d=1$ and consider the L\' evy process $L=(L(t))_{t\geq0}$.
We illustrate Theorem~\ref{wick} by taking
\[
Y(s)=L(s).
\]
By the Wick product calculus,
\[
\frac{d}{ds}L(s)^{\diamond 2}
=
2L(s)\diamond\overset{\bullet}{L}(s)
\]
in $(\mathcal S)^*$, and hence
\begin{equation}
\label{wick-square}
\int_0^t
L(s)\diamond\overset{\bullet}{L}(s)\,ds
=
\frac12L(t)^{\diamond2}.
\end{equation}

On the other hand, Theorem~\ref{wick} gives
\[
\int_0^t
L(s)\diamond\overset{\bullet}{L}(s)\,ds
=
\int_0^t L(s)L(ds),
\]
where the integral on the right-hand side is the Skorohod integral
defined through the chaos expansion.

Since $L(s)$ is adapted, in the one-parameter case the Skorohod
integral coincides with the It\^o integral. Thus
\[
\int_0^t L(s)L(ds)
=
\int_0^t L(s-)L(ds).
\]
Applying the It\^o formula for the square of a pure-jump L\' evy process,
we obtain
\begin{equation}
\label{ito-square}
L(t)^2
=
2\int_0^t L(s^-)L(ds)
+
\int_0^t\int_{\mathbb R_0}z^2N(ds,dz).
\end{equation}
Combining \eqref{wick-square} and \eqref{ito-square}, we obtain
\[
L(t)^{\diamond2}
=
L(t)^2
-
\int_0^t\int_{\mathbb R_0}z^2N(ds,dz).
\]
Equivalently,
\[
L(t)^{\diamond2}
=
L(t)^2-[L]_t,
\]
where $[L]_t$ denotes the quadratic variation of $L$.
\end{example}
\begin{example}
We illustrate Theorem \ref{wick} in the two-parameter case.
Let
\[
L(x)=L(x_1,x_2), \qquad x=(x_1,x_2)\in\mathbb{R}_+^2,
\]
be the L\' evy sheet and set
\[
Y(x)=L(x).
\]

Recall that the L\' evy sheet admits the representation
\[
L(x)
=
\int_{[0,x]}\int_{\mathbb{R}_0}
z\,\widetilde{N}(dy,dz),
\]
where
\[
[0,x]=[0,x_1]\times[0,x_2].
\]

Define the deterministic kernel
\[
\phi_x(y,z)
:=
\mathbf{1}_{[0,x]}(y)\,z,
\qquad
(y,z)\in\mathbb{R}^2\times\mathbb{R}_0.
\]
Consequently,
\[
Y(x)=L(x)=I_1(\phi_x).
\]
Hence, the process \(Y\) is Skorohod integrable
whenever the corresponding integrability condition in Definition
\ref{Skoro} is satisfied. Hence, by Theorem \ref{wick},
\[
\int_{[0,x]}L(y)\diamond\overset{\bullet} L(y)\,dy
=
\int_{[0,x]}L(y)\,L(dy),
\]
where the integral on the right-hand side is the Skorohod integral
defined by Definition \ref{Skoro}.

This identity should be distinguished from the multiparameter It\^o
integral. In particular, for \(d=2\), the Wick-calculus identity
\[
\frac12 L(x)^{\diamond2}
=
\int_{[0,x]}L(y)\diamond\overset{\bullet} L(y)\,dy
+
\int_{[0,x]}
\frac{\partial L}{\partial y_1}(y)
\diamond
\frac{\partial L}{\partial y_2}(y)\,dy
\]
shows that, in general,
\[
\frac12 L(x)^{\diamond2}
\neq
\int_{[0,x]}L(y)\,L(dy).
\]
There is therefore no contradiction with Theorem \ref{wick}: the
additional term is a genuinely multiparameter phenomenon and reflects
the distinction between the Skorohod integral and the multiparameter
It\^o integral. Moreover, to the best of our knowledge, a general
It\^o formula for multiparameter L\'evy sheets is not available.
\end{example}

\vspace{2mm} In view of the above theorem, the following terminology is natural.

\begin{definition}
\label{Wick-definition6.16} Suppose $Y$ is an $(\mathcal{S})^{\ast}$-valued
process such that
\[
\int_{\mathbb{R}^d}Y(x)\diamond\overset{\bullet}{L}(x)dx:\in(\mathcal{S})^{\ast
},
\]
then we call this integral \emph{the generalized Skorohod integral of $Y(x)$} with respect to $L(dx)$.
\end{definition}

Combining the properties above we get a powerful calculation technique for Skorohod 
integration. First of all, note that, by Theorem \ref{wick},
\begin{equation}
\int\limits_{0}^{x}\overset{\bullet}{L}(\xi)d\xi=L({x})\,. \label{3.63}%
\end{equation}
One can deduce that
\begin{equation}
\int\limits_{0}^{x}Z\diamond Y({\xi})\diamond\overset{\bullet}{L}({\xi})d\xi
=Z\diamond\int\limits_{0}^{x}Y({\xi})\diamond\overset{\bullet}{L}({\xi})d\xi,
\label{3.64}%
\end{equation}
if $Z$ does not depend on $\xi$. This is in contrast to the fact that for Skorohod
integrals we generally have
\begin{equation}
\int\limits_{0}^{x}Z\cdot Y({\xi})L(d\xi))\not =Z\cdot\int\limits_{0}%
^{x}Y({\xi})L(d\xi)\,, \label{3.65}%
\end{equation}
even if $Z$ does not depend on $\xi$.





\section{The population growth equation}
We illustrate the theory developed above by applying it to study the following model for population growth in a random environment driven by both multiparameter Brownian and  L\' evy white noise:

\begin{align}\label{heatgrowth}
 	 \frac{\partial \partial^{\alpha}}{\partial x\partial t^{\alpha}}Y(t,x)&= \mu(t,x) Y(t,x)+\sigma(t,x)Y(t,x) \diamond W(t,x)+\gamma(t,x)Y(t,x)\diamond V(t,x);\nonumber\\
     & (t,x)\in (0,\infty)\times \mathbb{R}^{d}.
 	 	\end{align}
with initial conditions

\begin{equation}
Y(0,x) = y_0, \quad Y(t,0) = y_0 \quad \text{(constant) for all } t,x \ge 0,
\end{equation}
where:
\begin{itemize}
    \item $\frac{\partial^\alpha}{\partial t^\alpha}$ is the Caputo fractional derivative of order $\alpha \in (0,1)$,
    \item $W(t,x) = \frac{\partial^2 B(t,x)}{\partial t \partial x}$ is the Brownian white noise (Brownian sheet),
    \item $V(t,x) = \frac{\partial^2 L(t,x)}{\partial t \partial x}$ is the L\'evy white noise (pure jump L\'evy sheet),
    \item $\mu(t,x)$, $\sigma(t,x)$, $\gamma(t,x)$ are deterministic functions,
    \item $y_0$ is a constant initial population density.
\end{itemize}
\begin{remark}
    Note that if $\alpha=1$ this corresponds to the following SDE driven by Brownian and L\' evy sheets:
\begin{align}\label{heatgrowth}
 	 Y(t,x)&= \int_0^t \int_0^x \mu(s,a) Y(s,a)ds da+\int_0^t \int_0^x \sigma(s,a)Y(s,a) B(ds,da)\nonumber\\
     &+\int_0^t \int_0^x \gamma(s,a)Y(s,a)L(ds,da);\nonumber\\
     & (t,x)\in (0,\infty)\times \mathbb{R}^{d}.
 	 	\end{align}
        \end{remark}
        
Before we state and prove our main result, we recall the following useful estimate on the norm of Wick products in the Kondratiev spaces:

\begin{proposition}(The V\r{a}ge inequality \cite{HOUZ}, Proposition 3.3.2)\\
Suppose $F\in (\S)_{-1,-k}$ and $G \in (\S)_{-1,-\ell}$ for some integers $k>\ell +1$.
Then
\begin{align}
 \|F \diamond G \|_{-1,-k} \leq A(k-\ell) \|F\|_{-1,-k} \|G\|_{-1,-\ell},  
\end{align}
where
\begin{align}
    A(k-\ell) = \sum_{\alpha} (2\mathbb{N})^{(\ell -k)\alpha}.
\end{align}
\end{proposition}
From this we deduce the following by induction on $n$:
\begin{corollary} \label{5.2}
    Suppose $K_i \in (\S)_{-1,-\ell}$ for some $\ell > 2; i=1,2, ...$
    Then for $k> \ell +1$ we have
    \begin{align}
        \|K_1 \diamond K_2 \diamond \cdots \diamond K_n\|_{-1,-k} \leq A(k-\ell)^{n-1} \Pi_{i=1}^n \|K_i\|_{-1,-\ell}
    \end{align}
    In particular, if $\|K_i\|_{-1,-\ell} \leq C$ for all $i$ then
  \begin{align}
        \|K_1 \diamond K_2 \diamond \cdots \diamond K_n\|_{-1,-k} \leq A(k-l)^{n-1} C^n.
    \end{align}  
\end{corollary}

We now state and prove our main result:
\begin{theorem}
    The solution $Y(t,x)$ of the fractional stochastic population growth equation \ref{heatgrowth} exists and is unique as an element of the Hida stochastic distribution space $(\S)^{*}$ and is explicitly represented as
\begin{align}
    Y(t,x)&=y_0 \sum_{n=0}^\infty \frac{1}{n!n!} \Big(\int_{[0,u]} K(v) dv \Big)^{\d n}
\end{align}
where \\
$K(v)=K(s,z)=\frac{1}{\Gamma(\alpha)}(t-s)^{\alpha -1}\Big(\mu(s,z) +\sigma(s,z)\overset{\bullet}{B}(s,z)  + \gamma(s,z) \overset{\bullet}{L}(s,z)\Big)$\\
and hence
\begin{align}
 &\int_{[0,u]} K(v) dv\nonumber\\
 &=\int_0^t \int_0^x \frac{1}{\Gamma(\alpha)}(t-s)^{\alpha -1}\Big(\mu(s,z) +\sigma(s,z)\overset{\bullet}{B}(s,z)  + \gamma(s,z) \overset{\bullet}{L}(s,z)\Big)ds dz\nonumber \\
 &= \int_0^t\int_0^x \frac{1}{\Gamma(\alpha)}(t-s)^{\alpha -1}\mu(s,z)ds dz\nonumber \\
 &+\int_0^t\int_0^x  \frac{1}{\Gamma(\alpha)}(t-s)^{\alpha -1}\sigma(s,z) B(ds,dz) +\int_0^t\int_0^x  \frac{1}{\Gamma(\alpha)}(t-s)^{\alpha -1}\gamma(s,z) L(ds,dz).
\end{align}    
\end{theorem}

\dproof
(i) We first prove existence and uniqueness:\\
By integrating  with respect to the spatial variable from $0$ to $x$:

\begin{align}
\frac{\partial^\alpha Y(t,x)}{\partial t^\alpha} - \frac{\partial^\alpha Y(t,0)}{\partial t^\alpha} &= \int_0^x \mu(t,z) Y(t,z) \, dz \notag \\
&\quad + \int_0^x \sigma(t,z) Y(t,z) \diamond W(t,z) \, dz \notag \\
&\quad + \int_0^x \gamma(t,z) Y(t,z) \diamond V(t,z) \, dz.
\end{align}

Since $Y(t,0) = y_0$, we obtain:

\begin{equation}\label{fracderiv}
\frac{\partial^\alpha Y(t,x)}{\partial t^\alpha} = \int_0^x \mu(t,z) Y(t,z) \, dz + \int_0^x \sigma(t,z) Y(t,z) \diamond W(t,z) \, dz + \int_0^x \gamma(t,z) Y(t,z) \diamond V(t,z) \, dz.
\end{equation}

Now apply the Caputo fractional integral of order $\alpha$ to both sides. Recall that the Caputo fractional integral $I^\alpha$ satisfies:
 \begin{equation}
I^\alpha \left( \frac{\partial^\alpha Y}{\partial t^\alpha} \right)(t,x) = Y(t,x) - Y(0,x).
\end{equation}
 
Hence:

\begin{align}
Y(t,x) &= y_0 + \frac{1}{\Gamma(\alpha)} \int_0^t (t-s)^{\alpha-1} \int_0^x \mu(s,z) Y(s,z) \, dz \, ds \notag \\
&\quad + \frac{1}{\Gamma(\alpha)} \int_0^t (t-s)^{\alpha-1} \int_0^x \sigma(s,z) Y(s,z) \diamond W(s,z) \, dz \, ds \notag \\
&\quad + \frac{1}{\Gamma(\alpha)} \int_0^t (t-s)^{\alpha-1} \int_0^x \gamma(s,z) Y(s,z) \diamond V(s,z) \, dz \, ds.
\end{align}   
That can be written as follow:

\begin{align}\label{integraleq}
Y(t,x) &= y_0 + \frac{1}{\Gamma(\alpha)} \int_0^t (t-s)^{\alpha-1} \int_0^x \mu(s,z) Y(s,z) \, dz \, ds \notag \\
&\quad + \frac{1}{\Gamma(\alpha)} \int_0^t (t-s)^{\alpha-1} \int_0^x \sigma(s,z) Y(s,z) \diamond B(ds, dz) \notag \\
&\quad + \frac{1}{\Gamma(\alpha)} \int_0^t (t-s)^{\alpha-1} \int_0^x \gamma(s,z) Y(s,z) \diamond L(ds, dz).
\end{align}


To solve this equation, we use the $S$-Transform method:


Let $F$ be a stochastic distribution in the Hida space. The S-transform of $F$ is defined by

\[
S(F)(\varphi) = \left\langle F , :e^{\langle \omega, \varphi \rangle}:\right\rangle.
\]

It satisfies

\[
S(F \diamond G)(\varphi) = S(F)(\varphi) S(G)(\varphi).
\]

We apply the S transform  to the equation \eqref{integraleq}:

Define
\[
u(t,x,\varphi) := S(Y(t,x))(\varphi).
\]

We have:

\[
S(\mu Y) = \mu(s,z)\ u(s,z,\varphi),
\]

\[
S(Y \diamond B(ds,dz)) = u(s,z,\varphi)\  \varphi(s,z)\ ds\ dz,
\]

\[
S(Y \diamond L(ds,dz)) = u(s,z,\varphi)\ \psi(s,z)\ ds\ dz,
\]

where
\[
\psi(s,z) = \int_{\mathbb{R}} (e^{\varphi(s,z,u)} - 1)\  \nu(du).
\]

Using all these, we obtain the following deterministic equation:

\begin{equation}
\begin{aligned}
u(t,x,\varphi)
&= y_0 \\
&+ \frac{1}{\Gamma(\alpha)} \int_0^t (t-s)^{\alpha-1}
\int_0^x \mu(s,z)\ u(s,z,\varphi)\  dz\  ds \\
&+ \frac{1}{\Gamma(\alpha)} \int_0^t (t-s)^{\alpha-1}
\int_0^x \sigma(s,z)\ u(s,z,\varphi)\ \varphi(s,z)\  dz\  ds \\
&+ \frac{1}{\Gamma(\alpha)} \int_0^t (t-s)^{\alpha-1}
\int_0^x \gamma(s,z)\ u(s,z,\varphi)\ \psi(s,z)\  dz\  ds.
\end{aligned}
\end{equation}

Define
\[
\beta_\varphi(s,z) = \mu(s,z) + \sigma(s,z)\ \varphi(s,z) + \gamma(s,z)\ \psi(s,z).
\]

Then
\begin{equation}\label{volt}
  u(t,x,\varphi) = y_0 + \frac{1}{\Gamma(\alpha)} \int_0^t (t-s)^{\alpha-1} \int_0^x \beta_\varphi(s,z)\ u(s,z,\varphi) dz  ds.  
\end{equation}


The S-transform reduces the stochastic equation to a deterministic linear fractional Volterra-type integral equation.

For a fixed test function $\varphi$, define the Volterra operator

\[
(K_\varphi f)(t,x)
=
\frac{1}{\Gamma(\alpha)}
\int_0^t (t-s)^{\alpha-1}
\int_0^x
\beta_\varphi(s,z)f(s,z)\,dz\,ds.
\]
Then equation \eqref{volt} can be written as

\[
u=y_0+K_\varphi u.
\]

We define the Picard sequence by
\[
u_0(t,x,\varphi)=y_0,
\]

and

\[
u_{n+1}(t,x,\varphi)
=
y_0+K_\varphi u_n(t,x,\varphi).
\]
The first iterates are

\[
u_1
=
y_0+K_\varphi y_0,
\]

\[
u_2
=
y_0+K_\varphi y_0+K_\varphi^2y_0,
\]

\[
u_3
=
y_0+K_\varphi y_0+K_\varphi^2y_0+K_\varphi^3y_0.
\]
By induction,

\[
u_n(t,x,\varphi)
=
\sum_{k=0}^{n}
K_\varphi^k y_0.
\]
We now compute $K_\varphi^k y_0$.\\ For $k=1$,

\[
K_\varphi y_0
=
\frac{y_0}{\Gamma(\alpha)}
\int_0^t
(t-s_1)^{\alpha-1}
\int_0^x
\beta_\varphi(s_1,z_1)\,dz_1\,ds_1.
\]
 
 For $k=2$,

\[
K_\varphi^2 y_0
=
\frac{y_0}{\Gamma(\alpha)^2}
\int_{0<s_2<s_1<t}
(t-s_1)^{\alpha-1}
(s_1-s_2)^{\alpha-1}
\]
\[
\qquad\qquad\times
\int_{0<z_2<z_1<x}
\beta_\varphi(s_1,z_1)
\beta_\varphi(s_2,z_2)
\,dz_2\,dz_1\,ds_2\,ds_1.
\]
Proceeding recursively, one obtains

\[
K_\varphi^k y_0
=
\frac{y_0}{\Gamma(\alpha)^k}
\int_{0<s_k<\cdots<s_1<t}
\prod_{j=1}^{k}
(s_{j-1}-s_j)^{\alpha-1}
\]

\[
\qquad\qquad\times
\int_{0<z_k<\cdots<z_1<x}
\prod_{j=1}^{k}
\beta_\varphi(s_j,z_j)
\,dz_k\cdots dz_1\,
ds_k\cdots ds_1,
\]

where $s_0=t$.
Assume that for every fixed $\varphi$ and every compact $[0,T]\times[0,X]$,
\[
B_\varphi(T,X)
:=
\sup_{(t,x)\in[0,T]\times[0,X]}
|\beta_\varphi(t,x)| < \infty.
\]
Then for any bounded function $f$,
\begin{align}
|K_\varphi f(t,x)|
&\le
\frac{B_\varphi(T,X)}{\Gamma(\alpha)}
\int_0^t (t-s)^{\alpha-1}
\int_0^x |f(s,z,\varphi)|\,dz\,ds.
\end{align}
Using the definition of the iterated integral in $\mathbb{R}^m$,
\[
\int_0^x dz = \prod_{i=1}^m x_i,
\quad \text{when } f \equiv 1,
\]
we obtain the bound
\begin{align}
\int_0^x |f(s,z)|\,dz
\le
\Big(\prod_{i=1}^m x_i\Big)\|f\|_\infty.
\end{align}
Let
\[
X=(X_1,\ldots,X_m)\in\mathbb{R}_+^m,
\qquad
[0,X]:=\prod_{i=1}^m[0,X_i].
\]
For \(x=(x_1,\ldots,x_m)\in[0,X]\), we have
\[
0\leq x_i\leq X_i,\qquad i=1,\ldots,m.
\]
Hence\begin{align}
\|K_\varphi f\|_\infty
\le
\frac{B_\varphi(T,X)}{\Gamma(\alpha)}
\Big(\prod_{i=1}^m X_i\Big)
\int_0^t (t-s)^{\alpha-1} ds\;\|f\|_\infty.
\end{align}
 Since
\[
\int_0^t (t-s)^{\alpha-1} ds
=
\frac{t^\alpha}{\alpha}
=
\frac{t^\alpha}{\Gamma(\alpha+1)}\Gamma(\alpha),
\]
we get
\begin{align}
\|K_\varphi f\|_\infty
\le
C_{\varphi,T,X}\, t^\alpha \|f\|_\infty,
\end{align}
where
\[
C_{\varphi,T,X}
=
\frac{B_\varphi(T,X)\prod_{i=1}^m X_i}{\Gamma(\alpha+1)}.
\]
By induction,
\[
\|K_\varphi^k y_0\|_\infty
\le
|y_0|
\frac{\big(C_{\varphi,T,X} T^\alpha\big)^k}{\Gamma(k\alpha+1)}.
\]
Since $\Gamma(k\alpha+1)$ grows super-exponentially, the series
\[
u(t,x,\varphi)
=
\sum_{k=0}^{\infty} K_\varphi^k y_0
\]
converges uniformly on $[0,T]\times[0,X]$.
Using linearity of $K_\varphi$,
\[
K_\varphi u
=
\sum_{k=1}^{\infty} K_\varphi^k y_0,
\]
hence
\[
y_0 + K_\varphi u
=
\sum_{k=0}^{\infty} K_\varphi^k y_0
=
u.
\]
For each fixed $\varphi \in \mathcal{S}(\mathbb{R}^d)$, there exists a unique solution
\[
u(t,x,\varphi)
\in C([0,T]\times[0,X])
\]
given by the Picard-Neumann series
\begin{equation}
    u(t,x,\varphi)
=
\sum_{k=0}^{\infty} K_\varphi^k y_0.
\end{equation}

\
It remains to reconstruct the stochastic solution from its
$S$-transform. For each fixed $(t,x)$, consider the mapping
\[
\varphi \longmapsto u(t,x,\varphi).
\]

This mapping is a
\emph{$U$-functional}. More precisely, it verifies the following
\begin{enumerate}
\item for all $\varphi_1,\varphi_2 \in \mathcal{S}(\mathbb{R}^d)$, the function
\[
z \longmapsto u(t,x,z\varphi_1+\varphi_2),
\qquad z\in\mathbb{C},
\]
is entire;

\item there exist constants $C,a,p>0$ such that
\[
|u(t,x,z\varphi)|
\le
C\exp\!\big(a|z|^2\|\varphi\|_p^2\big),
\]
for all $z\in\mathbb{C}$ and $\varphi\in\mathcal{S}(\mathbb{R}^d)$.
\end{enumerate}

The characterization
theorem for the $S$-transform in the Hida-Kondratiev framework
 implies that there exists a unique
stochastic distribution
\[
Y(t,x)\in (\mathcal{S})^{*}
\]
such that
\[
S(Y(t,x))(\varphi)
=
u(t,x,\varphi),
\qquad
\varphi\in\mathcal{S}(\mathbb{R}^d).
\]

Therefore, $u(t,x,\varphi)$ is the $S$-transform of a unique stochastic
distribution $Y(t,x)$, and we define
\[
Y(t,x)
=
S^{-1}\big(u(t,x,\cdot)\big).
\]

Consequently, $Y(t,x)$ is the unique stochastic distribution solution
of the original stochastic fractional integral equation.\\

(ii)
Next, to get an explicit solution formula, we work directly in the Kondratiev spaces $(\S)_{-1,-r}$ using Wick calculus: \\

Note that the equation can be written
\begin{align}
Y(t,x)=y_0 + \int_0^t \int_0^x K(s,z) \diamond Y(s,z) ds dz, \label{4.1}
\end{align}
where
\begin{align}
K(s,z)=K_t(s,z)=\frac{1}{\Gamma(\alpha)}(t-s)^{\alpha -1}(\mu(s,z) +\sigma(s,z)\overset{\bullet}{B}(s,z)  + \gamma(s,z) \overset{\bullet}{L}(s,z)).\label{4.2}
\end{align}
Substituting for $Y(s,z)$ in \eqref{4.1} we get
\begin{align*}
Y(u)= y_0 + \int_0^{u} K(u_1) \d \big(y_0+ \int_0^{u_1}K(u_2)\d Y(u_2)du_2\big)du_1,
\end{align*}
where we put  $u_j=(s_j,z_j); j=1,2, ...$.\\ 
Repeating this we obtain
\begin{align}
&Y(u)=y_0\Big( 1 + \int_0^{u} K(u_1)du_1 + \int_0^{u} (\int_0^{u_1} K(u_1) \d K(u_2)du_2 ) du_1 \nonumber\\
& +  ... + \int_0^{u}  \int_0^{u_1} ... \int_0^{u_n} K(u_1)\d K(u_2) \d \cdots K(u_{n}) du_1 du_2 \cdots du_n) +R_{n} \label{4.5}\\
&=y_0
+y_0
\sum_{k=1}^{n}
\int_{\Delta_k(u)}
K(u_1)\diamond K(u_2)\diamond\cdots\diamond K(u_k)
\,du_k\cdots du_1
+
R_n(u)
\end{align}
where
\[
\Delta_k(u)
=
\{0<s_k<\cdots<s_1<t,\; 0<z_k<\cdots<z_1<x\}.
\]
and 
\begin{align}
R_{n}&= \int_0^{u} \big( \int_0^{u_1}\cdots \int_0^{u_n} K(u_1)\d K(u_2) \d ... K(u_{n+1}) \d Y(u_{n+1})du_1 du_2 ...du_n \big) du_{n+1}\\
&=\int_{\Delta_{n+1}(u)}
K(u_1)\diamond\cdots\diamond K(u_{n+1})
\diamond Y(u_{n+1})
\,du_{n+1}\cdots du_1.
\end{align}
By the V\r{a}ge inequality (Corollary \ref{5.2}), it follows that $R_{n} \rightarrow 0$ in $(\mathcal{S})_{-1,-r}$ as $n \rightarrow \infty$.

Define 
$$F(u_1,u_2,\cdots, u_n)=K(u_1)\diamond K(u_2) \diamond \cdots. \diamond K(u_n).$$
Then the function $F$ is symmetric in all its variables $u_1, u_2,\cdots, u_n$.
This is simply by the fact that $\diamond$ is a commutative multiplication.
\\

Therefore we get, by letting $n \rightarrow \infty$ in \eqref{4.5},
\begin{align}
&Y(t,x)= y_0 \sum_{n=0}^\infty \int_0^{u} \big( \int_0^{u_1} ... \int_0^{u_n} K(u_1)\d K(u_2) \d ... K(u_n) du_1 du_2 ...du_n\\
&=y_0
\sum_{n=0}^{\infty}
\int_{\Delta_n(u)}
K(u_1)\diamond K(u_2)\diamond\cdots\diamond K(u_{n})
\,du_{1}\cdots du_n\\
&=y_0 \sum_{n=0}^\infty \frac{1}{n! n!} \int_{[0,u]^n} K(u_1)\d K(u_2) \d ... \d K(u_n) du_1 ... du_n\\
&=y_0 \sum_{n=0}^\infty \frac{1}{n!n!} \Big(\int_{[0,u]} K(v) dv \Big)^{\d n} \label{Y}
\end{align}
where
\begin{align}
 &\int_{[0,u]} K(v) dv\nonumber\\
 &=\int_0^t \int_0^x \frac{1}{\Gamma(\alpha)}(t-s)^{\alpha -1}\Big(\mu(s,z) +\sigma(s,z)\overset{\bullet}{B}(s,z)  + \gamma(s,z) \overset{\bullet}{L}(s,z)\Big)ds dz\nonumber \\
 &= \int_0^t\int_0^x \frac{1}{\Gamma(\alpha)}(t-s)^{\alpha -1}\mu(s,z)ds dz\nonumber \\
 &+\int_0^t\int_0^x  \frac{1}{\Gamma(\alpha)}(t-s)^{\alpha -1}\sigma(s,z) B(ds,dz) +\int_0^t\int_0^x  \frac{1}{\Gamma(\alpha)}(t-s)^{\alpha -1}\gamma(s,z) L(ds,dz).
\end{align}
The first part of the proof shows that the solution is unique and belongs to  $(\S)^{*}$. Therefore we can conclude that the solution $Y(t,x)$ given by \eqref{Y} actually belongs to $(\S)^{*}$.
\fproof

\begin{remark}
    Observe that the solution obtained in \eqref{Y} is not the Wick exponential. Indeed, the Wick exponential is defined by
\[
\exp^{\diamond}(Z)
=
\sum_{n=0}^{\infty}
\frac{1}{n!}
Z^{\diamond n},
\]
whereas our solution is given by
\[
Y(t,x)
=
y_0
\sum_{n=0}^{\infty}
\frac{1}{(n!)^2}
Z(t,x)^{\diamond n},
\]
with
\[
Z(t,x)
=
\int_0^u K(v)\,dv.
\]
Recall that the modified Bessel function of the first kind of order zero is defined by the power series
\[
I_0(z)
=
\sum_{n=0}^{\infty}
\frac{1}{(n!)^2}
\left(\frac{z}{2}\right)^{2n},
\]
or equivalently,
\[
I_0(2\sqrt{z})
=
\sum_{n=0}^{\infty}
\frac{z^n}{(n!)^2}.
\]
Therefore, the series appearing in \eqref{Y} is precisely the power series of the modified Bessel function evaluated at the Wick variable \(Z\). This motivates the introduction of the Wick--Bessel function
\[
I_0^{\diamond}(2\sqrt{Z})
:=
\sum_{n=0}^{\infty}
\frac{Z^{\diamond n}}{(n!)^2},
\]
where the ordinary powers are replaced by Wick powers.
Consequently, the solution can be written formally as
\begin{equation}
    Y(t,x)
=
y_0
I_0^{\diamond}
\!\left(
2\sqrt{
\int_0^u K(v)\,dv
}
\right).
\end{equation}
\end{remark}

\section{Acknowledgments}
 This research was funded by the Research Council of Norway under project 363480 (NASTRAN-Talent).

\begin{thebibliography}{99}

\bibitem{B} Benth, F. E. (1993): Integrals in the Hida distribution space $(\S)^{*}$.
In T. Lindstr\o m, B. \O ksendal and A. S.Ustunel (editors): Stochastic Analysis and Related Topics. Gordon
and Breach, pp. 89-99.

\bibitem{S} Schoutens W. (2000). \textit{Stochastic Processes and Orthogonal Polynomials}. Springer.

\bibitem{NS} Nualart, D., \& Schoutens, W. (2000). Chaotic and predictable representations for L\' evy processes. \textit{Stochastic Processes and their Applications}.

 
\bibitem{AMSW} Adler, R.J., Monrad, D., Scissors, R.H. \& Wilson, R. (1983): Representations, decompositions and sample
function continuity of random fields with independent increments. Stoch. Process. Appl. 15(1),3-30.
  
\bibitem{CMV} Capelas de Oliveira, E.; Mainardi \& F.Vaz, J. (2011): Models based on Mittag Leffler functions for anomalous relaxation in dielectrics. ariXir: 1106.1761  v2[cond-mat.Stat-mech] 13 Feb. 2014.

\bibitem{DW} Dalang, R. C. \& Walsh, J. B. (1992): The sharp Markov property of L\'evy sheets. The Annals of Probability 20 (2), 591-626.

\bibitem{DOP1} Di Nunno, G., \O ksendal, B. \& Proske, F. (2004):
White noise analysis for L\' evy processes. 
J. Functional Analysis 206 (2004),109 -148.

  \bibitem{DOP} Di Nunno, G., \O ksendal, B. \& Proske, F. (2009): Malliavin Calculus for L\' evy Processes with Applications to Finance. Springer.
  
 \bibitem{DMO} Draouil, O., Moulay-Hachemi, Y. \& \O ksendal, B. (2025): The stochastic heat inclusion with fractional time driven by  time-space Brownian noise. Submitted.
 	arXiv:2506.08275
  \bibitem{fedotov2007}
Fedotov, S. and Iomin, A. (2007).
  Migration and proliferation dichotomy in tumor cell invasion.
 Physical Review Letters, 98(11):118101.  
    
 \bibitem{GV} Gel'fand, I. M. \& Vilenkin, N. Ya.(1971): Generalized Functions. Vol. 4. Academic Press [Harcourt Brace Jovanovich Publishers], New York.
 
\bibitem{H} Hida, T. (1980): Brownian Motion. Springer, New York.

\bibitem{HKPS} Hida, T., Kuo, H-H., Potthoff, J. \& Streit, L. (1993): White Noise: An Infinite Dimensional Calculus. Springer, New York.

   \bibitem{Hu} Hu, Y. (2019): Some recent progress on stochastic heat equations. Acta Mathematica Scientia 39B(3); 874-914.
   
\bibitem{HOUZ}
Holden, H., \O ksendal,B., Ub\o e, J. \& Zhang, T. (2009)  Stochastic Partial Differential Equations, Springer.Second Edition.

 \bibitem{IR} Iafrate, F. \& Ricciutu, C.(2024): Some families of random fields related to multiparameter L\' evy processes. J. Theoretical Probability 37: 3055-3088. https://doi.org/10.1007/s10959-024-01351-3.
 
 \bibitem{Kilbas} Kilbas, A.A., Srivastava,H.M. \& Trujillo, J. J. (2006): 
	 {\emph{ Theory and Applications of Fractional Differential Equations }}, Elsevier  Science B.V, Amsterdam. 
     
\bibitem{KKS} Kochubel, A. N., Kondratiev, Y. \& da Silva, J. L. (2021): On fractional heat equation. Fractional Calculus \& Applied Analysis 24 (1), 73-87.

\bibitem{LOU} Lindstr\o m, T. , \O ksendal, B. and Ub\o e, J.  (1992): Wick multiplication and Ito-
Skorohod stochastic differential equations. In S. Albeverio, J. E. Fenstad, H. Holden
and T. Lindstr\o m (editors): Ideas and Methods in Mathematical Analysis, Stochastics,
and Applications. Cambridge Univ. Press, pp. 183-206.

   \bibitem{MS} Meerschaert, Mark M., Sikorskii, Alla (2019): Stochastic Models for Fractional Calculus, $2^{nd}$ edition. De Greuter.
   
   \bibitem{ML} Mittag-Leffler, M.G. (1903): Sur la nouvelle fonction E(x). \emph{C. R. Acad. Sci. Paris} {\bf 137} 554--558.
   
  \bibitem{MO1} Moulay Hachemi, R.Y.  \& \O ksendal, B. (2023): The fractional stochastic heat equation
driven by time-space white noise. Fractional Calculus \& Applied Analysis 26, 513-532. Springer.
https://doi.org/10.1007/s13540-023-00134-7.

\bibitem{murray2003}
Murray, J.D. (2003).
 Mathematical Biology II: Spatial Models and Biomedical Applications (3rd ed.). Springer-Verlag.
\bibitem{NualartSchoutens} Nualart, D. \& Schoutens, W. (2000): Chaotic and predictable representations for L\'evy processes.
Stochastic Process. Appl., 90(1):109-122.
\bibitem{OP} \O ksendal, B. \& Proske, F. (2004):
White noise of Poisson random measures. 
Potential Analysis 21 (2004), 375-403.
\bibitem{PK} Patel, A. \& Kosko, B. (2007): Levy Noise Benefits in Neural Signal Detection, 2007 IEEE International Conference on Acoustics, Speech and Signal Processing - ICASSP '07, Honolulu, HI, USA, 2007, pp. III-1413-III-1416, doi: 10.1109/ICASSP.2007.367111.  
\end{thebibliography}



\end{document}
@@@@@@@@@@@@@@@@@@@@@@@@@@@@@
\begin{theorem}
\label{th3.11} Choose $\varphi \in L^2(\mathbb{R}^d)$ and define the random variable 
$w_{\varphi}$ by $w_{\varphi}(\omega)=\langle \omega, \varphi\rangle; \varphi \in L^2(\R^d)$. Then 
\begin{equation}
w_{\varphi} ^{\diamond n}= ||\varphi ||^{n} h_n(\tfrac{w_{\varphi}}{|| \varphi||}) \label{herm} 
\end{equation}
where $||\varphi|| = ||\varphi||_{L^2(\mathbb{R}^d) }$.\newline

\end{theorem}

This result can be used to get an explicit formula for the Wick exponential:

\begin{theorem}{\bf The Wick exponential}

Let $\phi\in L^2(\R^d).$ Then
\begin{align*}
\exp^{\diamond}\Big(\int_{\RR^d} \phi(x) B(dx)\Big)=\exp\Big( \int_{\RR^d} \phi(x) B(dx)-\tfrac{1}{2} \int_{\RR^d} \phi^2(x)dx\Big).
\end{align*}
\end{theorem}

\dproof
We may assume that
$\phi=c\eta_1$, in which case we get, with $w(\phi)=\int_{\RR^d} \phi(x) B(dx)$,
\begin{align*}
\exp^\diamond[w(\phi)]&=\sum_{n=0}^\infty\frac{1}{n!}w(\phi)^{
\diamond
n}=\sum_{n=0}^\infty\frac{1}{n!}c^n\langle\omega,\eta_1\rangle^{\diamond n}\\
&=\sum_{n=0}^\infty\frac{c^n}{n!}H_{\epsilon_1}^{\diamond
n}(\omega)=\sum_{n=0}^\infty\frac{c^n}{n!}H_{n\epsilon_1}(\omega)\\
&=\sum_{n=0}^\infty\frac{c^n}{n!}h_n(\langle\omega,\eta_1\rangle)=\exp\Big[c
\langle\omega,
\eta_1\rangle-\tfrac{1}{2}c^2\Big]\\
&=\exp\bigg[w(\phi)-\tfrac{1}{2}\Vert\phi\Vert^2\bigg],
\end{align*}
where we have used the generating property of the Hermite polynomials.
\fproof
\subsection{Wick multiplication and It\^{o} integration}

One of the most
striking features of the Wick product is its relation to
It\^{o} integration. In short, this relation can be
expressed as
\begin{theorem}
Let $Y(x)$ be an It\^ o integrable process. Then 
\begin{align}\int\limits_{\R^d}Y(x)
B(dx)=\int\limits_{\R^d}Y(x)\diamond W(x)dx. \label{wis}
\end{align}
\end{theorem}
Here the left hand side denotes the It\^o integral of the stochastic process
$Y(x)=Y(x_1,x_2, ... , x_d,\omega)$ with respect to $B(dx)=B(dx_1 dx_2 ... dx_d)$, while the right hand side is to be
interpreted as an $({\cal S})^\ast$-valued (Pettis) Lebesgue integral of $Y \d W$  in $(\mathcal{S})^*$.

\vskip0.3cm
This relation explains why the Wick product is so natural and important
in stochastic calculus. It is also the key to the fact that It\^{o} calculus (with
It\^{o}'s formula, etc.) with ordinary multiplication is equivalent to ordinary
calculus with Wick multiplication. To illustrate this,
consider the example with $d=1, x=t$ and $Y(t)=B(t)\cdot\chi_{_{[0,T]}}(t)$.
Then by the It\^{o} formula the left hand side of \eqref{wis} becomes

\begin{align}\int\limits_0^T
B(t)dB(t)=\frac{1}{2}B^2(T)-\frac{1}{2}T\quad\text{(assuming}\;B(0)=0),\label{3.12} 
\end{align}
while (formal) Wick calculation makes the right hand side equal to
\begin{align*}
\int\limits_0^T B(t)\diamond W(t)dt=\int\limits_0^T B(t)\diamond
B'(t)dt=\frac{1}{2}B(T)^{\diamond 2},
\end{align*}
which is equal to \eqref{3.12}  in virtue of \eqref{herm}.











%%%%%%%%%%%%%%%%%%%%%%%%%%%%%%%%%%%%%%%%%%%%%
\subsection*{How can our model and the solution help?}

Our mathematical research offers a useful basis for clinical decision-making with precise implications for patient care. While $Y(t,x)$ is interpreted as a  tumor cell density (or its logarithm), this translates into clinical utility in several key areas.\\
At the first step, doctors need traditional imaging modalities such as MRI, CT scans that give a primitive hypothesis of the actual state of tumor. Our model uses this information as the initial condition $Y(0,x)$ to generate a patient-specific prediction of invasion cells.

The model can predict microscopic extensions beyond the visible tumor boundary on conventional imaging, also helping surgeons to decide about optimal resection margins while preserving critical areas.
%%%%%%%%%%%%%%%%%%%%%%%%%%%%%%%%%%%%%%%
The multi-parameter component of 
 L\'evy noise $\gamma V(t,x)$ captures sudden and discontinuous changes in tumor behavior that are clinically significant. Unexpected jumps in $Y(t,x)$ despite ongoing treatment can lead to the emergence of new tumor masses. Separating the continuous evolution from the jump events helps distinguish treatment-related inflammation from actual tumor growth.
 \section{Time-space white noise}
In this subsection we summarize the construction of time-space white noise and its properties. A general reference for this section is Chapter 2 in \cite{HOUZ}.
\vskip 0.3cm
Let $n$ be a fixed natural number. Later we will set $n= 1 + d$.
Define $\Omega={\cal S}'(\RR^n)$, equipped
with the weak$^{*}$ topology. This space will be the base of our basic
probability space, which we explain in the following:
\vskip 0.3cm
As events we will use the family $\mathcal{F}={\cal
B}({\cal S}'(\RR^n))$ of Borel subsets of ${\cal S}'(\RR^d)$, and our
probability measure $\P$ is defined by the following result: 
\begin{theorem}{\bf (The Bochner--Minlos theorem)}\\
There exists a unique probability measure $\P$ on ${\cal B}({\cal S}'(\RR^n))$
with the following property:
$$\E[e^{i\langle\cdot,\phi\rangle}]:=\int\limits_{\cal S'}e^{i\langle\omega,
\phi\rangle}d\mathbb{P}(\omega)=e^{-\tfrac{1}{2} \Vert\phi\Vert^2};\quad i=\sqrt{-1}$$
for all $\phi\in{\cal S}(\RR^n)$, where
$\Vert\phi\Vert^2=\Vert\phi\Vert^2_{L^2(\RR^n)},\quad\langle\omega,\phi\rangle=
\omega(\phi)$ is the action of $\omega\in{\cal S}'(\RR^n)$ on
$\phi\in{\cal S}(\RR^n)$ and $\E=\E_{\P}$ denotes the expectation
with respect to  $\P$.
\end{theorem} 
We will call the triplet $({\cal S}'(\RR^n),{\cal
B}({\cal S}'(\RR^n)),\P)$ the {\it  white noise probability
space\/}, and $\P$ is called the {\it white noise probability measure}.
The measure $\P$ is also often called the (normalised) {\it Gaussian
measure\/} on ${\cal S}'(\RR^n)$. It is not difficult to prove that if $\phi\in L^2(\RR^n)$ and
we choose $\phi_k\in{\cal S}(\RR^n)$ such that $\phi_k\to\phi$ in $L^2(\RR^n)$,
then
$$\langle\omega,\phi\rangle:=\lim\limits_{k\to\infty}\langle\omega,\phi_k\rangle
\quad\text{exists in}\quad L^2(\P)$$
and is independent of the choice of $\{\phi_k\}$. In
particular, if we define
$$\widetilde{B}(x):=\widetilde{B}(x_1,\cdots,x_n,\omega)=\langle\omega,\chi_
{[0,x_1]\times\cdots\times[0,x_n]}\rangle; \quad  
x=(x_1,\cdots,x_n)\in\RR^n,$$
where $[0,x_i]$ is interpreted as $[x_i,0]$ if $x_i<0$,
then $\widetilde{B}(x,\omega)$ has an $x$-continuous version $B(x,\omega)$, which
becomes an \emph{$n$-parameter Brownian sheet}, in the following sense:
 By an \emph{$n$-parameter Brownian sheet} we mean a family
$\{B(x,\cdot)\}_{x\in\RR^n}$ of random variables on a probability space
$(\Omega,{\cal F},\P)$ such that
\begin{itemize} 
\item
$B(0,\cdot)=0\quad\text{almost surely with respect to } \P,$
\item
$\{B(x,\omega)\}$ is a continuous and Gaussian stochastic process 
\item
For all $x=(x_1,\cdots,x_n)$,
$y=(y_1,\cdots,y_n)\in\RR_+^n$,
$B(x,\cdot),\,B(y,\cdot)$ have the covariance
$\prod_{i=1}^n x_i\wedge y_i$. For general $x,y\in\RR^n$ the covariance is
$\prod_{i=1}^n\int_\RR\theta_{x_i}(s)\theta_{y_i}(s)ds$, where
$\theta_x(t_1,\dots,t_n)=\theta_{x_1}(t_1)\cdots\theta_{x_n}(t_n)$, with
\begin{equation*} 
\theta_{x_j}(s)=
\begin{cases}
1\quad \text{ if } 0<s\leq x_j\\
-1\quad \text{ if } x_j<s\leq 0\\
0 \quad \text{ otherwise}
\end{cases}
\end{equation*} 
\end{itemize} 
It can be proved that the process $\widetilde{B}(x,\omega)$ defined above has a modification $B(x,\omega)$ which satisfies all these properties.
This process $B(x,\omega)$ then becomes an \emph{$n$-parameter Brownian sheet}. 
\vskip0.3cm
We remark that for $n=1$ we get the classical (1-parameter) Brownian motion
$B(t)$ if we restrict ourselves to $t\geq 0$. 
\vskip0.3cm
With this definition of the Brownian sheet it is natural to define the
$n$-parameter Wiener--It\^{o} integral of $\phi\in L^2(\RR^n)$ by
$$\int\limits_{\RR^n}\phi(x)dB(x,\omega):=\langle\omega,\phi\rangle;\quad
\omega\in{\cal S}'(\RR^d).$$
We see that by using the Bochner--Minlos theorem we have obtained an easy construction of
$n$-parameter Brownian motion/sheet that works for any parameter dimension $n$. Moreover, we get a representation of the space $\Omega$ as the dual $\mathcal{S}'(\R^d)$ of the Fr\' echet space $\mathcal{S}(\R^d)$. This is an advantage in many situations, for example in the construction of the Hida-Malliavin derivative, which can be regarded as a stochastic gradient on $\Omega$.
See e.g. \cite{DOP} and the references therein.\\
In the following we put $n=1+d$ and let 
 $$B(t,x)=B(t,x,\omega); t \geq 0, x \in \R^d, \omega \in \Omega$$ 
 denote the (1-dimensional) time-space Brownian sheet with probability law $\P$. 
\subsection{Chaos expansion in terms of Hermite polynomials}

The {\it Hermite polynomials\/} $h_n(x); x \in \R$ are defined by
$$h_n(x)=(-1)^n
e^{{1}/{2}x^2}\frac{d^n}{dx^n}(e^{-{1}/{2}x^2});\quad
n=0,1,2,\cdots.\leqno{(2.2.1)}$$
Thus the first Hermite polynomials are
\begin{align}
h_0(x)&=1,\;h_1(x)=x,\;h_2(x)=x^2-1,\;h_3(x)=x^3-3x,\cr
h_4(x)&=x^4-6x^2+3,\;h_5(x)=x^5-10x^3+15x,\cdots.
\end{align}
The \emph{Hermite functions} $\xi_n(x)$ are defined by
$$\xi_n(x)=\pi^{-{1}/{4}}((n-1)!)^{-{1}/{2}}e^{-{1}/{2}x^2}h_{n-1}
(\sqrt{2}x);\quad n=1,2,\cdots.$$
Some useful properties of $h_n$ and $\xi_n$ are:
\begin{itemize}
\item
$\xi_n\in{\cal S}(\R)\quad\text{for all}\;n.$
\item
 The collection $\{\xi_n\}^\infty_{n=1}$ constitutes an
orthonormal basis for $L^2(\R).$
\item
$\sup\limits_{x\in\R}|\xi_n(x)|=O(n^{-{1}/{12}}).$
\end{itemize}
 Proofs of these statements can be found in Hille and
Phillips (1957), Chapter 21.
\vskip 0.2cm
We will use these functions to define an orthogonal basis for $L^2(\P)$.  In the following, we let
$\beta=(\beta_1,\cdots,\beta_n)$ denote $n$-dimensional multi-indices with
$\beta_1,\cdots,\beta_n\in\N$. By the above it follows that the family of tensor
products
$$\xi_\beta:=\xi_{(\beta_1,\cdots,\beta_n)}:=\xi_{\beta_1}\otimes\cdots
\otimes
\xi_{\beta_n}\;;\;\beta\in\N^n$$
forms an orthonormal basis for $L^2(\R^n)$. Let
$\beta^{(j)}=(\beta_1^{(j)},\beta_2^{(j)},\cdots,\beta_n^{(j)})$ be
the $j$th multi-index number in some fixed ordering of all $n$-dimensional
multi-indices $\beta=(\beta_1,\cdots,\beta_n)\in\N^n$. We 
assume that this ordering has the property that
$$i<j\Rightarrow\beta_1^{(i)}+\beta_2^{(i)}+\cdots+\beta_n^{(i)}\leq
\beta_1^{(j)}+\beta_2^{(j)}+\cdots+\beta_n^{(j)},$$
i.e., that the $\{\beta^{(j)}\}^\infty_{j=1}$ occur in increasing order.


Now define
$$\eta_j:=\xi_{\beta^{(j)}}=\xi_{\beta_1^{(j)}}\otimes\cdots\otimes\xi_{
\beta_n
^{(j)}};\;j=1,2,\cdots.$$
We will need to consider multi-indices of arbitrary length. To
simplify the notation, we regard multi-indices as elements of the space
$(\N_0^{\N})_c$ of all sequences $\alpha=(\alpha_1,\alpha_2,\cdots)$
with elements $\alpha_i\in\N_0$ and with compact support,
i.e., with only finitely many $\alpha_i\not=0$. We write
$${\cal J}=(\N_0^{\N})_c.$$

\begin{definition}
Let $\alpha=(\alpha_1,\alpha_2,\cdots)\in{\cal J}$. Then we define
$$H_\alpha(\omega):=\prod\limits_{i=1}^\infty
h_{\alpha_i}(\langle\omega,\eta_i\rangle);\quad\omega\in{\cal
S}'(\R^n).\leqno{(2.2.10)}$$
\end{definition}

\begin{theorem}(Wiener--It\^{o} chaos expansion theorem) 
Every
$f\in L^2(\P)$ has a unique representation
$$f(\omega)=\sum\limits_{\alpha\in{\cal J}}c_\alpha
H_\alpha(\omega)$$
where $c_\alpha\in\R$ for all $\alpha$.
\end{theorem}

Moreover, we have the isometry
$$\Vert
f\Vert^2_{{L}^2(\P)}=\sum\limits_{\alpha\in{\cal
J}}\alpha!c_\alpha^2.$$

\begin{example}
$i)$ The 1-dimensional smoothed white noise $w(\phi,\omega):=\langle\omega,\phi\rangle$
 has the expansion
\begin{align}
    w(\phi,\omega)&=\langle\omega,\phi\rangle=\langle\omega,\sum
\limits_{j=1}^\infty(\phi,\eta_j)\eta_j\rangle\cr
&=\sum\limits_{j=1}^\infty(\phi,\eta_j)\langle\omega,\eta_j\rangle=\sum
\limits_{j=1}^\infty(\phi,\eta_j)H_{\epsilon^{(j)}}(\omega),
\end{align}
where  $\epsilon^{(j)}=(0,0,\cdots,1,\cdots)$ with 1 on entry number $j$, 0 otherwise. The
convergence is in $L^2(\P)$.
\end{example}

In other words,
$$w(\phi,\omega)=\sum\limits_\alpha c_\alpha H_\alpha(\omega)$$
with
$c_\alpha=(\phi,\eta_j)$  if  $\alpha=\epsilon^{(j)};$
\quad 0 otherwise.
\vskip0.3cm
$ii)$ The 1-dimensional, $n$-parameter Brownian motion
$B(x,\omega)$ is defined by
$$B(x,\omega)=\langle\omega,\psi\rangle,$$
where
$$\psi(y)=\chi_{[0,x_1]\times\cdots\times[0,x_n]}(y).$$
Proceeding as above, we write
$$\psi(y)=\sum\limits_{j=1}^\infty(\psi,\eta_j)\eta_j(y)=\sum\limits_{j=1}^
\infty\int\limits_0^x\eta_j(u)du\,\eta_j(y),$$
where we have used the multi-index notation
$$\int\limits_0^x\eta_j(u)du=\int\limits_0^{x_n}\cdots\int\limits_0^{x_1}
\eta_j(u_1,\cdots,u_n)du_1\cdots
du_n=\prod\limits_{k=1}^n\int\limits_0^{x_k}\xi_{\beta_k^{(j)}}(t_k)
dt_k$$
when $x=(x_1,\cdots,x_n)$. Therefore,
$$B(x,\omega)=\langle\omega,\sum\limits_{j=1}^\infty\ (\int\limits_0^x\eta_j
(u)du) \eta_j\rangle=\sum\limits_{j=1}^\infty(\int\limits_0^x\eta_j(u)du)
\langle\omega,\eta_j\rangle,$$
so $B(x,\omega)$ has the expansion
$$B(x,\omega)=\sum\limits_{j=1}^\infty (\int\limits_0^x\eta_j(u)du)\,H_
{\epsilon^{(j)}}(\omega).$$

\begin{example}(Brownian white noise)
Since the Brownian sheet $B(t,x)$ is $(t,x)$-continuous a.s., we can for a.a. $\omega \in \Omega$ define its derivatives with respect to $t$ and $x$ in the sense of a tempered distribution, i.e. in $\S'(\R^n)=\S'(\R \times \R^d)$.
 Thus we define the time-space white noise $W(t,x)=W(t,x,\omega)$ by
 8
\begin{equation}
  	W(t,x)=\frac{\partial}{\partial t}\frac{\partial^{d}B(t,x)}{\partial x_{1}...\partial x_{d}}. \label{WN}
  \end{equation}
In particular, for $d=1,$ and putting $x_1=t,$ we get the familiar identity
$$W(t)=\frac{d}{dt}B(t)\hbox{ in } \S^{'}(\R).$$ 

However, we can also formally obtain a chaos expansion of $W(t,x)$ by differentiating the expansion for $B(t,x)$:
\begin{align}
W(t,x)&=\frac{\partial}{\partial t}\frac{\partial^{d}}{\partial x_{1}...\partial x_{n}}B(t,x)=\frac{\partial}{\partial t}\frac{\partial^{d}}{\partial x_{1}...\partial x_{d}}\sum\limits_{j=1}^\infty (\int_0^t \int\limits_0^x\eta_j(s,v)ds dv)\,H_{\epsilon^{(j)}}(\omega)\nonumber\\
&=\sum\limits_{j=1}^\infty \eta_j(t,x)\,H_{\epsilon^{(j)}}(\omega).
\end{align}
We shall see that this expansion defines $W(t,x)$ as an element in a stochastic distribution space, namely the Hida space, defined as follows:
\end{example}

\begin{definition}
\label{Def10.4}{\bf (The Hida spaces)} 
\begin{itemize}
\item[(i)] \quad\index{L\'evy-Hida st8ochastic test functin space}
{\bf Stochastic test functions $(\mathcal{S})$.} Let \ $(\mathcal{S})$ consist of all $%
\varphi =\sum_{\alpha \in \mathcal{J}}a_{\alpha }H_{\alpha }\in L^{2}(\P)$
such that%
\begin{equation}
\left\Vert \varphi \right\Vert _{k}^{2}:=\sum_{\alpha \in \mathcal{J}%
}a_{\alpha }^{2}\alpha !(2\mathbb{N})^{k\alpha }<\infty \text{ for \emph{all}
}k\in \mathbb{N},  \label{10.18}
\end{equation}%
equipped with the projective topology, where%
\begin{equation}
(2\mathbb{N})^{k\alpha }=\prod_{j\geq 1}(2j)^{k\alpha _{j}},
\end{equation}%
if $\alpha =(\alpha _{1,}\alpha _{2,},...)\in \mathcal{J}$.
\item[(ii)]\quad
{\bf Stochastic distributions $(\mathcal{S})^{* }$ .} \index{L\'evy-Hida stochastic distribution space}
Let $(\mathcal{S})^{* }$ consist of all expansions $F=\sum_{\alpha \in \mathcal{%
J}}b_{\alpha }H_{\alpha }$ such that%
\begin{equation}
\left\Vert F\right\Vert _{-q}^{2}:=\sum_{\alpha \in \mathcal{J}}b_{\alpha
}^{2}\alpha !(2\mathbb{N})^{-q\alpha }<\infty \text{ for \emph{some} }q\in
\mathbb{N}.  \label{norm}
\8end{equation}%
endowed with the inductive topology.
The space $(\mathcal{S})^{* }$ is the
dual of $(\mathcal{S}).$ If $F=\sum_{\alpha \in \mathcal{J}}b_{\alpha
}H_{\alpha }\in (\mathcal{S})^{* }$ and $\varphi =\sum_{\alpha \in
\mathcal{J}}a_{\alpha }H_{\alpha }\in (\mathcal{S}),$ then the action of $F$
on $\varphi $ is%
\begin{equation}
\left\langle F,\varphi \right\rangle =\sum_{\alpha \in \mathcal{J}}a_{\alpha
}b_{\alpha }\alpha !.  \label{10.21}
\end{equation}
\item[(iii)]\quad
{\bf Generalized expectation.}\index{generalized expectation}
If  $F= \sum_{\alpha \in \mathcal{J}} a_{\alpha }H_{\alpha } \in (\mathcal{S})^{* }$, we define
the generalized expectation $\E[F]$ of $F$ by
$$
\E[F] = a_0.
$$ 
Note that $\E[H_{\alpha }] = 0$ for all $\alpha \ne 0$.
Therefore the generalized expectation coincides with the usual expectation if $F \in L^2(\P)$. 
\end{itemize}
\end{definition}
Using this definition one can prove that $W(t,x)$ can be regarded as a Hida stochastic distribution:
\begin{proposition}
8$$W(t,x,\omega)=\sum\limits_{j=1}^\infty \eta_j(t,x)\,H_{\epsilon^{(j)}}(\omega)\in(\S(\R \times \R^d))^{*} \quad\text{for each} \quad t,x\in \R \times \R^d.$$
\end{proposition}
We refer to \cite{HOUZ}, Chapter 2 for proofs and more information.
%%%%%%%%%%%%%%%%%%%%%%%%%%%%%%%%%
%%%%%%%%%%%%%%%%%%%%%%%%%%%%%%%%
\textcolor{blue}{
We define $\F_{t,x}$ to be the filtration generated by the Brownian sheet $B(s,a)$ and L\' evy sheet $L(r,b)$ for $s,r\leq t,a,b \leq x$.\\
Let $(\S)^{*}=(\S(\R_{+} \times \R^d))^{*}$ denote the set of Hida stochastic distributions. Let $Y(t,x)=Y(t,x,\omega); (t,x) \in \R_{+} \times \R^d, \omega \in \Omega$ be an $(\S)^{*}$-valued process and let $G: (\S)^{*} \mapsto \mathcal{B}((\S)^{*})$ be a set-valued map.}

%%%%%%%%%%%%%%%%%%%%%%%%%%%%%%
%%%%%%%%%%%%%%%%%%%%%%%%%%%%%%%
\subsection{Mathematical background}
 We now explain our model more precisely:\\
 

%, where  $\mathcal{B}(L_a^{\infty}(\P))$ denotes the family of all Borel subsets of the set $L_a^{\infty}(\P)$ of all bounded $\F_{t,x}$-adapted %processes with compact support.



Let $n$ be a given natural number. Let $\mathcal{S}=\mathcal{S}(\mathbb{R}^n)$\label{simb-028} be the
space of rapidly decreasing smooth real
functions $f$
on $\mathbb{R}^n$
equipped with the family of seminorms:\label{simb-029} 
\begin{equation*}
\Vert f \Vert_{k,\gamma} := \sup_{x \in \mathbb{R}^n}\big\{ (1+|x|^k) \vert
\partial^ \gamma f(x)\vert \big\}< \infty,
\end{equation*}
where $k = 0,1,...$, $\gamma=(\gamma_1,...,\gamma_n)$ is a multi-index with $%
\gamma_j= 0,1,...$ $(j=1,...,n)$ and\label{simb-030} 
\begin{equation*}
\partial^\gamma f (x):= \frac{\partial^{|\gamma|}}{\partial
x_1^{\gamma_1}\cdots \partial x_n^{\gamma_n}}f(x)
\end{equation*}
for $|\gamma|=\gamma_1+ ... +\gamma_n$.


Then
$\mathcal{S}=\mathcal{S}(\mathbb{R}^n)$ is a
Fr\'echet space.

Let $\mathcal{S}^{\prime }=\mathcal{S}^{\prime }(\mathbb{R}^{n})$\label%
{simb-031} be its dual, called the space of \emph{tempered distributions}. 
\index{tempered distributions} Let $\mathcal{B}$ denote the family of all
Borel subsets of $\mathcal{S}^{\prime }(\mathbb{R}^{n})$ equipped with the
weak* topology. If $\Phi \in \mathcal{S}^{\prime }$ and $f \in \mathcal{%
S}$ we let \label{simb-033} 
\begin{equation}
\Phi (f) \text{ or } \langle \Phi ,f \rangle  \label{3.1}
\end{equation}%
denote the action of $\Phi$ on $f$.
\end{example}

\begin{example}
\begin{itemize}
\item
{(Evaluations)}
For $y \in \R$ define the function $\delta_y$ on $\S(\R)$ by $\delta_y(\phi)=\phi(y)$. Then $\delta_y$ is a tempered distribution.\vskip 0.2cm
\item
{(Derivatives)} Consider the function $D$, defined for $\phi \in \S(\mathbb{R})$ by $D[\phi]=\phi^{\prime}(y)$. Then  $D$ is a tempered distribution. \vskip 0.2cm
\item
{(Distributional derivative)}\\
 Let $T$ be a tempered distribution, i.e. $T \in \S^{'}(\mathbb{R}) $. We define the distributional derivative $T^{'}$ of $T$ by
 $$ T^{'}[\phi]=-T[\phi^{'}]; \quad \phi \in \S.$$
 Then $T^{'}$ is again a tempered distribution.
 \end{itemize}
 \end{example}


 In the following we will apply this to the case when $n=1+d$ and $y=(t,x) \in \R \times \R^d$.


\subsection*{Mild solutions}

A solution $y(t,x)$ is called \emph{mild} if $\E|Y^2(t,x)] < \infty$ for all $t,x$.\\
In this case the (a-priori) distribution valued solution can be represented as a pointwise-defined physical quantity.

It is easy to verify that the results in \cite{MO1} cally over to the include the L\' evy noise case, and we can conclude:\\
 - In the classical case where $\alpha = 1$, mild solution exists in one spatial dimension $d = 1$.\\

 - In the case where $\alpha > 1$, mild solutions exist for $d=1$ or $d=2$. 

 - In the sub-diffusion case where $\alpha < 1$, there are no mild solutions for any spatial dimension.
We now pass from abstract existence to a rigorous chaos representation
in the case $\sigma=0$, where the system is driven only by a pure jump
L\' evy sheet.

Since $u(t,x,\cdot)$ is a $U$-functional on the nuclear space
$\mathcal{S}(\mathbb{R}^d)$, it is analytic in every direction and
satisfies exponential growth estimates. Hence, it admits the Taylor
expansion at $0$:
\[
u(t,x,\varphi)
=
\sum_{n=0}^{\infty}
\frac{1}{n!}
D^n u(t,x,0)[\varphi^{\otimes n}],
\qquad \varphi\in\mathcal{S}(\mathbb{R}^d).
\]

By the nuclearity of $\mathcal{S}(\mathbb{R}^d)$ and the Schwartz kernel
theorem, each continuous symmetric $n$-linear form
$D^n u(t,x,0)$ can be uniquely represented by a distribution
\[
F_n(t,x)\in\mathcal{S}'((\mathbb{R}^d)^n)
\]
such that
\[
D^n u(t,x,0)[\varphi_1,\dots,\varphi_n]
=
\langle F_n(t,x), \varphi_1\otimes\cdots\otimes \varphi_n\rangle.
\]
In particular,
\[
D^n u(t,x,0)[\varphi^{\otimes n}]
=
\langle F_n(t,x), \varphi^{\otimes n}\rangle.
\]

Thus,
\[
u(t,x,\varphi)
=
\sum_{n=0}^{\infty}
\langle F_n(t,x), \varphi^{\otimes n}\rangle.
\]

Recall that in the present setting ($\sigma=0$), the stochastic noise
is a pure jump L\' evy sheet represented by a compensated Poisson random
measure $\widetilde N(d\xi,dz)$ on $\mathbb{R}_+^d \times \mathbb{R}_0$.
The $S$-transform inversion corresponds therefore to multiple stochastic
integrals with respect to this Poisson random measure, namely
\[
S^{-1}(\langle F_n, \varphi^{\otimes n}\rangle)
=
I_n^{\widetilde N}(F_n),
\]
where $I_n^{\widetilde N}(F_n)$ denotes the $n$-th Hida multiple
stochastic integral associated with the L\' evy sheet.

Using linearity of $S^{-1}$ and the fact that the series converges in
$U(\mathcal{S})$ the space of $U$-functionals, we obtain
\[
Y(t,x)
=
S^{-1}(u(t,x,\cdot))
=
\sum_{n=0}^{\infty}
I_n^{\widetilde N}(F_n(t,x)).
\]

By the characterization theorem for $U$-functionals
(Holden--\O ksendal-Ub\o e-Zhang), there exists
$p\ge0$ such that the chaos coefficients associated with
$u(t,x,\cdot)$ satisfy the Kondratiev summability condition.
Consequently,
\[
Y(t,x)\in(\mathcal S)_{-p}^{-1}
\subset (\mathcal S)^*.
\]


Therefore, the stochastic distribution solution admits the rigorous
Poisson-type chaos expansion
\[
Y(t,x)
=
\sum_{n=0}^{\infty}
I_n^{\widetilde N}(F_n(t,x))
\in (\mathcal{S})^*,
\]
which is the unique solution of the stochastic fractional integral
equation driven by a L\' evy sheet.
\textcolor{red}{Observe that, although
\[
|\Delta_n(u)|
=
\frac{1}{(n!)^2}|[0,u]^n|,
\]
one cannot replace the integration over $\Delta_n(u)$ by an integration
over the whole box $[0,u]^n$ multiplied by the factor $(n!)^{-2}$.
Indeed, such a replacement is valid only when the integrand is symmetric
with respect to permutations of the variables
$u_1,\dots,u_n$.
In the present situation, the integrand
\[
K(u,u_1)\diamond K(u_1,u_2)\diamond\cdots\diamond K(u_{n-1},u_n)
\]
is not symmetric. For example, when $n=2$ we have
\[
K(u,u_1)\diamond K(u_1,u_2),
\]
where
\[
K(u,u_1)
=
\frac{(t-s_1)^{\alpha-1}}{\Gamma(\alpha)}
\Big(
\alpha(s_1,z_1)
+\sigma(s_1,z_1)\dot B(s_1,z_1)
+\gamma(s_1,z_1)\dot L(s_1,z_1)
\Big),
\]
and
\[
K(u_1,u_2)
=
\frac{(s_1-s_2)^{\alpha-1}}{\Gamma(\alpha)}
\Big(
\alpha(s_2,z_2)
+\sigma(s_2,z_2)\dot B(s_2,z_2)
+\gamma(s_2,z_2)\dot L(s_2,z_2)
\Big).
\]
Interchanging $u_1$ and $u_2$ yields
\[
K(u,u_2)\diamond K(u_2,u_1),
\]
which is generally different from
\[
K(u,u_1)\diamond K(u_1,u_2),
\]
since the factors
\[
(t-s_1)^{\alpha-1}, \qquad (s_1-s_2)^{\alpha-1}
\]
are replaced by
\[
(t-s_2)^{\alpha-1}, \qquad (s_2-s_1)^{\alpha-1}.
\]
Therefore the integrand depends on the ordering of the variables
$u_1,\dots,u_n$ and is not permutation invariant. Consequently,
the identity
\[
\int_{\Delta_n(u)}
K(u,u_1)\diamond\cdots\diamond K(u_{n-1},u_n)
\,du_n\cdots du_1
=
\frac{1}{(n!)^2}
\int_{[0,u]^n}
K(u,u_1)\diamond\cdots\diamond K(u_{n-1},u_n)
\,du_n\cdots du_1
\]
is, in general, false.}
We now pass from abstract existence to a rigorous chaos representation
in the case $\sigma=0$, where the system is driven only by a pure jump
L\'evy sheet.

In the Hida-Kondratiev framework for L\' evy white noise, the $S$-transform
is defined with respect to the Poisson-Wick structure. In particular, for a test function
$\varphi \in \mathcal{S}(\mathbb{R}^d \times \mathbb{R}_0)$, we set
\[
\psi(s,z)
=
\int_{\mathbb{R}_0}
\big(e^{\varphi(s,z,u)} - 1\big)\,\nu(du).
\]

Since $u(t,x,\cdot)$ is a $U$-functional on the nuclear L\' evy test space,
it is analytic in every direction and satisfies exponential growth estimates.
Hence, it admits the Taylor expansion at $0$:
\[
u(t,x,\varphi)
=
\sum_{n=0}^{\infty}
\frac{1}{n!}
D^n u(t,x,0)[\varphi^{\otimes n}],
\qquad \varphi \in \mathcal{S}(\mathbb{R}^d \times \mathbb{R}_0).
\]

By the nuclearity of the test space and the Schwartz kernel theorem,
each continuous symmetric $n$-linear form $D^n u(t,x,0)$ can be uniquely represented by a distribution
\[
F_n(t,x) \in \mathcal{S}'\big((\mathbb{R}^d \times \mathbb{R}_0)^n\big)
\]
such that
\[
D^n u(t,x,0)[\varphi_1,\dots,\varphi_n]
=
\langle F_n(t,x), \varphi_1 \otimes \cdots \otimes \varphi_n \rangle.
\]

Thus,
\[
u(t,x,\varphi)
=
\sum_{n=0}^{\infty}
\langle F_n(t,x), \varphi^{\otimes n} \rangle.
\]

In the present L\' evy setting, the correct chaos structure is not Gaussian
Wiener chaos but Poisso-Wick (Charlier-type) chaos generated by the
compensated Poisson random measure $\widetilde N(d\xi,dz,du)$ associated with the L\' evy sheet.

The inverse $S$-transform acts on Poisson-Wick monomials as
\[
S^{-1}\!\left(\langle F_n, \psi^{\otimes n} \rangle\right)
=
I_n^{\widetilde N}(F_n),
\]
where $I_n^{\widetilde N}(F_n)$ denotes the $n$-th multiple stochastic integral
with respect to the compensated Poisson random measure $\widetilde N$.

Using linearity of $S^{-1}$ and the fact that $u(t,x,\cdot)$ is a $U$-functional in the L\' evy-Hida sense,
we obtain
\[
Y(t,x)
=
S^{-1}(u(t,x,\cdot))
=
\sum_{n=0}^{\infty}
I_n^{\widetilde N}(F_n(t,x)).
\]

By the characterization theorem for $U$-functionals
(Holden-Oksendal-UbOe-Zhang), there exists $p \ge 0$ such that the chaos
coefficients satisfy the Kondratiev growth condition. Consequently,
\[
Y(t,x) \in (\mathcal{S})_{-p}^{-1} \subset (\mathcal{S})^*.
\]

Therefore, the stochastic distribution solution admits the rigorous
Poisson-Wick chaos expansion
\[
Y(t,x)
=
\sum_{n=0}^{\infty}
I_n^{\widetilde N}(F_n(t,x))
\in (\mathcal{S})^*,
\]
which is the unique solution of the stochastic fractional integral equation
driven by a pure jump L\' evy sheet.

Now we will give the chaos expansion in the pure jump L\' evy case ($\sigma=0$):

We consider the deterministic $S$-transform equation
\[
u(t,x,\varphi)
=
y_0 + (K_\varphi u)(t,x),
\]
where
\[
(K_\varphi f)(t,x)
=
\frac{1}{\Gamma(\alpha)}
\int_0^t (t-s)^{\alpha-1}
\int_{\mathbb{R}^d}
\beta_\varphi(s,z)\, f(s,z)\,dz\,ds.
\]

The L\' evy structure enters through
\[
\beta_\varphi(s,z)
=
\alpha(s,z)
+
\gamma(s,z)\,\psi(\varphi)(s,z),
\]
where the L\' evy functional is
\[
\psi(\varphi)(s,z)
=
\int_{\mathbb{R}_0}
\big(e^{\varphi(s,z,u)}-1\big)\nu(du).
\]

By the Picard iteration

\[
u_0(t,x)=y_0,
\qquad
u_{n+1}(t,x)=y_0 + K_\varphi u_n(t,x),
\]
so that
\[
u(t,x,\varphi)
=
\sum_{n=0}^{\infty} K_\varphi^n y_0.
\]
By induction,
\[
K_\varphi^n y_0
=
\frac{y_0}{\Gamma(\alpha)^n}
\int_{\Delta_n(t,x)}
\prod_{j=1}^n (s_{j-1}-s_j)^{\alpha-1}
\prod_{j=1}^n \beta_\varphi(s_j,z_j)\,
d\mathbf{s}\,d\mathbf{z},
\]
where $s_0=t$ and
\[
\Delta_n(t,x)=\{0<s_n<\cdots<s_1<t\}.
\]
We separate the drift and jump parts by the following:
\[
\beta_\varphi(s,z)
=
\alpha(s,z)
+
\gamma(s,z)\psi(\varphi)(s,z).
\]

Hence,
\[
\prod_{j=1}^n \beta_\varphi(s_j,z_j)
=
\sum_{k=0}^n
\sum_{\substack{A\subset\{1,\dots,n\}\\ |A|=k}}
\left(\prod_{j\in A}\gamma(s_j,z_j)\psi(\varphi)(s_j,z_j)\right)
\left(\prod_{j\notin A}\alpha(s_j,z_j)\right).
\]
The inverse $S$-transform acts on L\' evy functionals according to the
Poisson--Wick rule:
\[
S^{-1}\big(\psi(\varphi)(s,z)\big)
=
\int_{\mathbb{R}_0} u\,\widetilde N(ds,dz,du),
\]
and more generally,
\[
S^{-1}\Big(
\psi(\varphi)(s_1,z_1)\cdots \psi(\varphi)(s_k,z_k)
\Big)
=
I_k^{\widetilde N}(s_1,z_1,\dots,s_k,z_k),
\]
where $I_k^{\widetilde N}$ denotes the $k$-th multiple stochastic integral
with respect to the compensated Poisson random measure $\widetilde N$.

We use the expansion
\[
u(t,x,\varphi)
=
\sum_{n=0}^{\infty}
\langle F_n(t,x), \varphi^{\otimes n}\rangle.
\]

Thus, $F_n(t,x)$ is obtained by collecting all terms of total L\' evy order $n$,
i.e. all contributions containing exactly $n$ factors of $\psi(\varphi)$.

We write
\[
F_n(t,x)
=
\frac{y_0}{\Gamma(\alpha)^n}
\sum_{k=0}^n
F_{n,k}(t,x),
\]
where $k$ counts the number of jump contributions,
with
\[
F_{n,k}(t,x)
=
\int_{\Delta_n(t,x)}
\prod_{j=1}^n (s_{j-1}-s_j)^{\alpha-1}
\,K_{n,k}(s_1,z_1,\dots,s_n,z_n)\,
d\mathbf{s}\,d\mathbf{z},
\]
where
\[
K_{n,k}
=
\sum_{\substack{A\subset\{1,\dots,n\}\\ |A|=k}}
\left(\prod_{j\in A}\gamma(s_j,z_j)\right)
\left(\prod_{j\notin A}\alpha(s_j,z_j)\right)
\cdot \mathcal{I}_{A}^{\widetilde N},
\]
and
\[
\mathcal{I}_{A}^{\widetilde N}
:=
S^{-1}\Big(
\prod_{j\in A}\psi(\varphi)(s_j,z_j)
\Big)
=
I_{|A|}^{\widetilde N}.
\]

Applying the inverse $S$-transform yields the L\' evy chaos expansion
\[
Y(t,x)
=
S^{-1}(u(t,x,\cdot))
=
\sum_{n=0}^{\infty}
I_n^{\widetilde N}(F_n(t,x)).
\]

Moreover,
\[
Y(t,x)\in(\mathcal{S})_{-p}^{-1}\subset(\mathcal{S})^*,
\]
so that $Y(t,x)$ is the unique stochastic distribution solution
of the original fractional L\' evy-driven Volterra equation.
%%%%%%%%%%%%%%%%%%%%%%%%%%%%%%%%%%%%%%%%%%%%%
%%%%%%%%%%%%%%%%%%%%%%%%%%%%%%%%%%%%%%%%%%%
\section{ The stochastic heat equation with fractional time
driven by time-space Brownian and L\' evy white noise}
As a first application of the multiparameter L\' evy white noise theory, we consider the fractional heat equation driven by time-space Brownian and L\' evy white noise. Our equation generalizes the fractional stochastic heat equation studied in \cite{MO1}, since here we account not only for the time-space Brownian noise but also for the L\' evy white noise: 
\begin{equation}\label{heat1}
 	 \frac{\partial^{\alpha}}{\partial t^{\alpha}}Y(t,x)= \l \Delta Y(t,x)+\sigma W(t,x)+\gamma V(t,x);\; (t,x)\in (0,\infty)\times \mathbb{R}^{d}.
 	 	\end{equation}
       Here $d\in\mathbb{N}=\{1,2,...\}$ is the space dimension, $x \in \R^d$ denotes a point in space, $t\geq 0$ denotes the time,  $\frac{\partial^{\alpha}}{\partial t^{\alpha}}$ is the Caputo derivative of order $\alpha \in (0,2)$, and $\l>0$, $\sigma, \gamma \in \mathbb{R}$ are given constants,
   
  	\begin{equation}
  		\Delta Y =\sum_{j=1}^{d}\frac{\partial^{2}Y}{\partial x_{j}^{2}}(t,x)
  	\end{equation}
  is the Laplacian operator with respect to $x$, and
 
  \begin{equation}
  	W(t,x)=W(t,x,\omega)=\frac{\partial}{\partial t}\frac{\partial^{d}B(t,x)}{\partial x_{1}...\partial x_{d}} 
  \end{equation}
is time-space white noise, where $$B(t,x)=B(t,x,\omega); t\geq 0, x=(x_1, x_2, ... ,x_d) \in \R^d, \omega \in \Omega$$
 is time-space Brownian sheet with probability law $\P$, and

 \begin{align}
     V(t,x) = V(t,x,{\omega)=\frac{\partial}{\partial t}\frac{\partial^{d}L(t,x)}{\partial x_{1}...\partial x_{d}} }
 \end{align}
 is a L\'evy white noise, where $L(t,x)$ is a $1+d$ - parameter L\'evy sheet.
The boundary conditions are
\begin{align}
    Y(0,x)&=\delta(x)\text{ (the Dirac measure at  } x), \label{1.4}\\
    \lim_{x \rightarrow +/- \text{ }\infty}Y(t,x)&=0.\label{1.5}
\end{align}

\begin{remark}
    $W(t,x)$ and $V(t,x)$ model two different types of noise, which both can be represented as elements of $(\S)^{*}$: \\
    W (t, x) is a continuous Gaussian noise, while $V(t,x)$ represents noise coming from jumps from a L\' evy sheet. \\
    By construction we can assume that $W(t,x)$ and $V(t,x)$ are independent under the same probability measure $\P$ on $\Omega=\S'(\R_+ \times \R^d)$.
\end{remark}
By extending the method in \cite{MO1} to our equation, we obtain the following result, which we state without proof: 
\begin{theorem} \label{th1}
	The unique solution $Y(t,x) \in (\S)^{*}$ of the fractional stochastic heat equation \eqref{heat1} - \eqref{1.5} is given by	
	\begin{align}\label{th2}
		Y(t,x)&=I_1 + I_2+I_3,
	\end{align}
	where
 		\begin{align}
			I_1=(2\pi)^{-d} \int_{\mathbb{R}^d} e^{ixy} E_{\alpha}(- \l t^{\alpha} |y|^2) dy
			=(2\pi)^{-d} \int_{\mathbb{R}^d} e^{ixy}\sum_{k=0}^{\infty} \frac{(- \l t^{\alpha} |y|^2)^k}{\Gamma(\alpha k +1)}dy, 
		\end{align}
	and
			\begin{align}
			I_2=\sigma (2\pi)^{-d} \int_{0}^{t}(t-r)^{\alpha -1}\int_{\mathbb{R}^{d}}\left(\int_{\mathbb{R}^{d}}e^{i(x-z)y}\sum_{k=0}^{\infty}\frac{(-\l (t-r)^{\alpha}|y|^2)^{k}}{\Gamma(\alpha k+\alpha))}dy\right) B(dr,dz) , 
		\end{align}
        and
        \begin{align}
			I_3
			=\sigma (2\pi)^{-d} \int_{0}^{t}(t-r)^{\alpha -1}\int_{\mathbb{R}^{d}}\left(\int_{\mathbb{R}^{d}}e^{i(x-z)y}\sum_{k=0}^{\infty}\frac{(-\l (t-r)^{\alpha}|y|^2)^{k}}{\Gamma(\alpha k+\alpha))}dy\right) L(dr,dz). 
		\end{align}
 
\end{theorem}
%%%%%%%%%%%%%%%%%%%%%%%%%%%%%%%%%%%%%%%%%%%%%%
$$$$$$$$$$$$$$$$$$$$$$$$$$$$$$$$$$
$$$$$$$$$$$$$$$$$$
We proceed to establish some useful properties of generalized Skorohod integrals.

\begin{lemma}
\label{Wick-lemma6.16} Suppose  $G(x)\in(\mathcal{S}%
)_{1,-q}$ for all $x\in\mathbb{R}^d$, for some $q\in\mathbb{N}_{+}$. 
Then
\begin{equation}
   \int_{\mathbb{R}^d} \vert< G(x)\diamond\overset{\bullet}{L}(x),f > \vert dx
\leq\Vert f \Vert_{1,q}\Big( \int_{\mathbb{R}^d} \Vert G(x) \Vert^{2}_{1,-q}dx
\Big)^{1/2}
\end{equation}
for all $f \in (\S)_{1,q}.$
\end{lemma}

\dproof
Choose $\rho \in [-1,1]$ and suppose $G(x)= \sum_{\alpha\in \mathcal{J}}a_{\alpha}(x)K_{\alpha} \in (\S)_{\rho,q}$. Choose 
$f= \sum_{\beta\in \mathcal{J}}b_{\beta}K_{\beta} \in (\S)_{-\rho,-q}$.
Then

\begin{align}
<G(x) \diamond \overset{\bullet}{L}(x),f>
&=m_2 <\sum_{\alpha,k} a_{\alpha}(x) e_{k}(x)K_{\alpha+\epsilon^{(k,1)}}, \sum_{\beta\in \mathcal{J}}b_{\beta}K_{\beta}>\\
&=m_2 \sum_{\alpha,k} a_{\alpha}(x) e_{k}(x) b_{\alpha+\epsilon^{(k,1)}} (\alpha+\epsilon^{(k,1)})!\,.
\end{align}

Hence 
\begin{align}
\int_{\mathbb{R}^d}&\vert <G(x) \diamond \overset{\bullet}{L}(x),f> \vert dx
\leq  m_2 \sum_{\alpha,k} \vert b_{\alpha+\epsilon^{(k,1)}}\vert \alpha! (\alpha_{k} +1) \int_{\mathbb{R}^d} \vert a_{\alpha}(x)e_{k}(x)\vert dx\\
\leq& m_2 \sum_{\alpha,k} \vert b_{\alpha+\epsilon^{(k,1)}} \vert \alpha! (\alpha_{k} +1) \Big( \int_{\mathbb{R}^d} a_{\alpha}^{2}(x)dx \Big)^{1/2}\\
\leq &m_2 \Big( \sum_{\alpha,k} b_{\alpha+\epsilon^{(k,1)}}^{2}  \Big((\alpha+\epsilon^{(k,1)})!\Big)^{1+\rho} (2\mathbb{N})^{q(\alpha+\epsilon^{(k,1)})}  \Big)^{1/2} \\
&\cdot \Big( \sum_{\alpha,k} \Big( \int_{\mathbb{R}}a_{\alpha}^{2}(x)dx\Big)  \Big(\alpha!\Big)^{1-\rho} (\alpha_{k}+1)^{1-\rho} (2\mathbb{N})^{-q\alpha}(2k)^{-q}  \Big)^{1/2} \\
\leq& m_2 \Vert f\Vert_{\rho,q} \Big( \sum_{\alpha,k} \Big( \int_{\mathbb{R}^d}a_{\alpha}^{2}(x)dx\Big)     \Big(\alpha!\Big)^{1-\rho} (\alpha_{k}+1)^{1-\rho} (2k)^{-q} (2\mathbb{N})^{-q\alpha}  \Big)^{1/2} \\
\leq& m_2 \Vert f\Vert_{1,q} \Big( \int_{\mathbb{R}^d} \Vert G(x)\Vert^{2}_{1,-q}dx \Big)^{1/2}
\end{align}
if $\rho=1$.
\fproof


\vspace{2mm} \noindent Using this result we obtain the following theorem.

\begin{theorem}
\label{Wick-theorem6.17}

\begin{enumerate}
\item[(i)] Suppose $G:\mathbb{R}^d  \to (\mathcal{S})_{-1}$ satisfies
\[
\int_{\mathbb{R}^d} \Vert G(x) \Vert^{2}_{1,-q} dx < \infty, \quad\text{ for some }
q\in\mathbb{N}.
\]
Then
\[
\int_{\mathbb{R}^d} G(t)\diamond\overset{\bullet}{L}(x)dx \quad\text{ exists in }
(\mathcal{S})_{-1,-q}.
\]


\item[(ii)] Suppose $F(x)$, $F_{n}(x)$, $n=1,2,...$, are elements of
$(\mathcal{S})_{1,-q}$ for all $x\in\mathbb{R}^d$ and
\[
\int_{\mathbb{R}^d} \Vert F_{n}(x) - F(x) \Vert^{2}_{1,-q}dx \to 0, \qquad
n\to\infty.
\]
Then
\[
\int_{\mathbb{R}^d} F_{n}(x) \diamond\overset{\bullet}{L}(x)dx \to
\int_{\mathbb{R}^d} F (x) \diamond\overset{\bullet}{L}(x)dx, \quad n \to\infty,
\]
in $(\mathcal{S})_{-1,q}$.
\end{enumerate}
\end{theorem}

\dproof
(i) The proof follows from Lemma \ref{Wick-lemma6.16} and Definition \ref{definition3.13-added}:\\
The integral $\int_{\R^d} G(x) \diamond \overset{\bullet}{L}(x) dx $  exists as an element in the dual $(\S)_{-1,-q}$ because by Lemma \ref{Wick-lemma6.16} it is a bounded linear functional on $(\S)_{1,q}$.\\
(ii)\, By Lemma \ref{Wick-lemma6.16} we have
\begin{align*}
\vert < \int_{\R^d} \big( F_{n}(x)-F(x)\big)& \diamond \overset{\bullet}{L}(x)dx, f> \vert\\
\leq & \int_{\R^d} \vert < \big( F_{n}(x)-F(x)\big) \diamond \overset{\bullet}{L}(x), f> \vert dx\\
\leq &m_2 \Vert f \Vert_{1,q}\int_{\R^d} \Vert  F_{n}(x)-F(x) \Vert^{2}_{1,-q}dx \to 0, \quad n\to\infty.
\end{align*}
\fproof

\begin{theorem} \cite{LOU},\cite{B}.\label{3.5}
\label{wick} Assume that $Y(x)$, $x\in\mathbb{R}^d$, is a Skorohod integrable
stochastic process (recall Definition \ref{skoro}). Then $Y(x)
\diamond\overset{\bullet}{L}(x)$, $x\in\mathbb{R}^d$, is integrable in
$(\mathcal{S})^{*}$ and
\begin{equation}
\label{3.66}\int_{\mathbb{R}^d} Y(x) \diamond \overset{\bullet}{L}(x) dx= \int_{\mathbb{R}^d} Y(x) L(dx).
\end{equation}

\end{theorem}

\dproof
Assume that $Y(x)$ has the representation
$$
Y(x) = \sum_{\alpha\in \mathcal{J}} c_{\alpha}(x) K_{\alpha}
= \sum_{n=0}^{\infty} I_{n} \big( f_{n}(\cdot,x)\big).
$$
where
\begin{equation}
f_n(\cdot,x)=f_n (x^{(1)},z_{1};...; x^{(n)},z_{n};x)=\sum_{\left\vert \alpha \right\vert
=m}c_{\alpha }(x)\theta ^{\hat{\otimes }\alpha
}(x^{(1)},z_{1},..., x^{(n)},z_{n}).
\end{equation}
Then the left-hand side of \eqref{3.66} can be written 
\begin{equation}
\label{3.67}
\int_{\mathbb{R}^d} \big(\sum_{\alpha\in \mathcal{J}} c_{\alpha}(x) K_{\alpha}\big) \diamond
 \big(\sum_{k} e_{k}(x) K_{\varepsilon ^{(k,1)}}m_{2}\big) dx
=
\sum_{\alpha, k} (c_{\alpha},e_{k})_{L^{2}(\mathbb{R}^d)} K_{\alpha+\epsilon ^{(k,1)}}m_2.
\end{equation}
For the right-hand side we obtain, using that
$L(dx)=\int_{\R_0} z \widetilde{N}(dx,dz),$
\begin{align}
\int_{\mathbb{R}^d} Y(x) L(dx)\label{3.68}
& = \int_{\mathbb{R}^d} \sum_{n=0}^{\infty} I_{n} \big( f_{n}(\cdot,x)\big) L(dx)\nonumber\\
&= \sum_{n=0}^{\infty} \int_{\mathbb{R}^d} I_{n }\big( \sum_{|\alpha|=n}c_{\alpha}(x) \theta^{\hat\otimes\alpha}\big) L(dx) \\
& = \sum_{n=0}^{\infty} \int_{\mathbb{R}^d} I_{n }\big( \sum_{|\alpha|=n} \sum_{k=1}^{\infty}
(c_{\alpha}, e_{k})_{{}_{L^{2}(\mathbb{R}^d)}} e_{k}(x) \theta^{\hat\otimes\alpha} \big) L(dx)\\
& = \sum_{n=0}^{\infty} \int_{\mathbb{R}^d} \int_{\R_0}I_{n }\big( \sum_{|\alpha|=n} \sum_{k=1}^{\infty}
(c_{\alpha}, e_{k})_{{}_{L^{2}(\mathbb{R}^d)}} e_{k}(x) \theta^{\hat\otimes\alpha} \big) z\widetilde{N}(dx,dz)\\
& = m_2 \sum_{n=0}^{\infty} I_{n+1} \Big( \sum_{|\alpha|=n} \sum_{k=1}^{\infty}
(c_{\alpha}, e_{k})_{{}_{L^{2}(\mathbb{R}^d)}} \big( \theta^{\hat\otimes\alpha}\hat\otimes \theta_{k}\big) \Big)\\
&=m_2 \sum_{n=0}^{\infty} \sum_{|\alpha|=n} \sum_{k=1}^{\infty}
(c_{\alpha}, e_{k})_{{}_{L^{2}(\mathbb{R}^d)}} I_{n+1}\big( \theta^{\hat\otimes (\alpha + \epsilon^{(k,1)})}\big)\\
& =m_2 \sum_{\alpha, k}  (c_{\alpha},e_{k})_{{}_{L^{2}(\mathbb{R}^d )}} K_{\alpha+\epsilon ^{(k,1)}}.
\end{align}

Comparing \eqref{3.67} and \eqref{3.68}, we get \eqref{3.66}.
%\qedsymbol
\fproof

\section {Computation by using orthogonal polynomials}
In some cases one can compute these Wick powers by using that the Wick power of an integral with respect to $B$ is given by a Hermite polynomial $H_n$ and the Wick power of an integral with respect to $L$ is given by a Laguerre polynomial $\gamma_n$.\\

More precisely, if $\varphi$ is deterministic and $\| \varphi\|=1$, we have
\begin{align}
    (\int_{\R}\int_{\R^d} \varphi(s,z)B(ds, dz))^{\diamond n}=H_n(\int_{\R}\int_{\R^d} \varphi(s,z)B(ds,dz)).
\end{align}
Similarly, if $L$ is Gamma distributed and $\psi$ is deterministic and $\| \psi\|=1$, we would like to have
\begin{align}
    (\int_{\R}\int_{\R^d} \psi(s,z)L(ds, dz))^{\diamond n}=\gamma_n(\int_{\R}\int_{\R^d} \psi(s,z)L(ds,dz))
\end{align}

\subsection{Remark on Laguerre polynomial representations for Gamma noise}
It is important to distinguish between a Gamma distributed random variable and a
stochastic integral with respect to a Gamma process or Gamma sheet.
\\For a single Gamma random variable
\[
G\sim \mathrm{Gamma}(\alpha,\lambda),
\]
the orthogonal polynomial system associated with the Gamma measure is given by
the generalized Laguerre polynomials
\[
L_n^{(\alpha-1)}(\lambda G).
\]
The corresponding Gamma Wick (Appell) polynomials are obtained from the
generating function
\[
\sum_{n=0}^{\infty}\frac{t^n}{n!}P_n(G)
=
\exp(tG)
\left(1-\frac{t}{\lambda}\right)^{\alpha}.
\]
In this one-dimensional setting, the Wick powers can be represented as
\[
G^{\diamond n}=P_n(G),
\]
where \(P_n\) is a Laguerre-type polynomial. In particular, the second Wick
power has the form
\[
G^{\diamond2}=G^2-aG-b,
\]
where the constants \(a\) and \(b\) depend on the parameters
\(\alpha\) and \(\lambda\).
\\However, this representation does not extend directly to stochastic
integrals with respect to a Gamma process or a Gamma sheet. Indeed, let
\[
F=\int_0^t f(s)\Gamma(ds)
\]
be a first-order Gamma stochastic integral. The Wick square of \(F\) admits a
chaos decomposition containing higher-order and lower-order Gamma chaos terms,
and the first-order correction generally involves nonlinear functions of the
kernel, such as
\[
\int_0^t f(s)^2\Gamma(ds),
\]
which is not, in general, a deterministic multiple of \(F\).
Consequently, one cannot write in general
\[
F^{\diamond n}=L_n^{(\alpha-1)}(F)
\]
or any formula involving only the value of the random variable \(F\). The
Wick powers depend on the complete kernel defining the stochastic integral.\\The situation is therefore analogous to the Poisson case. For a single
Poisson random variable, the Wick powers are represented by Charlier
polynomials, whereas for a general Poisson stochastic integral one obtains
generalized Charlier chaos polynomials depending on the kernel.\\
Thus, for a Gamma sheet integral of the form
\[
G(t,x)=
\int_0^t\int_0^x
g(t,x;s,\xi)\Gamma(ds,d\xi),
\]
the correct object is a generalized Gamma chaos polynomial associated with
the kernel \(g_{t,x}\), and not a classical Laguerre polynomial evaluated at
the scalar random variable \(G(t,x)\).
\begin{proposition}
Let $T,X>0$ and let $0<\alpha<1$. Assume that there exists an integer
$\ell_0\geq 0$ such that
\begin{equation}
\sup_{(s,z)\in[0,T]\times[0,X]}
\|\mathcal K(s,z)\|_{-1,-\ell_0}
\leq C_{\mathcal K}<\infty.
\label{eq:K-bound}
\end{equation}
Let $r$ be an integer satisfying
\begin{equation}
r>\ell_0+1.
\label{eq:r-condition}
\end{equation}
Assume moreover that
\[
(s,z)\longmapsto \mathcal K(s,z)
\]
is continuous from $[0,T]\times[0,X]$ into
$(\mathcal S)_{-1,-\ell_0}$.

Then the series
\begin{equation}
Y(t,x)=y_0+\sum_{n=1}^{\infty}Y_n(t,x)
\label{eq:Y-series}
\end{equation}
converges uniformly on $[0,T]\times[0,X]$ with values in
$(\mathcal S)_{-1,-r}$ and
\begin{equation}
Y\in
C\big([0,T]\times[0,X];(\mathcal S)_{-1,-r}\big).
\end{equation}
Furthermore, if $R_n$ denotes the remainder obtained after $n$ iterations
of the Volterra equation, then
\begin{equation}
\sup_{(t,x)\in[0,T]\times[0,X]}
\|R_n(t,x)\|_{-1,-r}
\longrightarrow0,
\qquad n\to\infty.
\end{equation}
\end{proposition}

\begin{proof}
We consider the Volterra equation
\begin{equation}
Y(t,x)
=
y_0+
\int_0^t\int_0^x
K(t,x;s,z)\diamond Y(s,z)\,dz\,ds,
\label{eq:Volterra}
\end{equation}
where
\begin{equation}
K(t,x;s,z)
=
\frac{(t-s)^{\alpha-1}}{\Gamma(\alpha)}
\mathcal K(s,z),
\qquad 0<\alpha<1.
\label{eq:K-definition}
\end{equation}

We first introduce the iterated kernels. For convenience, put
\[
u_0=(t,x),\qquad
u_j=(s_j,z_j),\quad j\geq1,
\]
and set $s_0=t$. We define
\begin{equation}
K(u_{j-1};u_j)
=
\frac{(s_{j-1}-s_j)^{\alpha-1}}
{\Gamma(\alpha)}
\mathcal K(s_j,z_j).
\label{eq:iterated-K}
\end{equation}

For $n\geq1$, define
\begin{equation}
\Delta_n(t,x)
=
\left\{
0<s_n<\cdots<s_1<t,\quad
0<z_n<\cdots<z_1<x
\right\}.
\label{eq:Delta}
\end{equation}

The $n$-th term of the iterated series is
\begin{equation}
Y_n(t,x)
=
y_0
\int_{\Delta_n(t,x)}
K(u_0;u_1)\diamond\cdots\diamond
K(u_{n-1};u_n)
\,du_n\cdots du_1.
\label{eq:Yn}
\end{equation}

Notice that the terminal point changes at every iteration. In particular,
the product in \eqref{eq:Yn} is
\[
K(t,x;s_1,z_1)
\diamond
K(s_1,z_1;s_2,z_2)
\diamond\cdots\diamond
K(s_{n-1},z_{n-1};s_n,z_n),
\]
and not a product in which $(t,x)$ is kept fixed in every factor.

We now estimate the iterated Wick products.

By \eqref{eq:K-bound} and \eqref{eq:iterated-K},
\begin{equation}
\begin{aligned}
\|K(u_{j-1};u_j)\|_{-1,-\ell_0}
&=
\frac{(s_{j-1}-s_j)^{\alpha-1}}
{\Gamma(\alpha)}
\|\mathcal K(s_j,z_j)\|_{-1,-\ell_0}
\\
&\leq
C_{\mathcal K}
\frac{(s_{j-1}-s_j)^{\alpha-1}}
{\Gamma(\alpha)}.
\end{aligned}
\label{eq:K-factor-estimate}
\end{equation}

Recall the Våge inequality: if
$F\in(\mathcal S)_{-1,-r}$ and
$G\in(\mathcal S)_{-1,-\ell_0}$ with
$r>\ell_0+1$, then
\begin{equation}
\|F\diamond G\|_{-1,-r}
\leq
A(r-\ell_0)
\|F\|_{-1,-r}
\|G\|_{-1,-\ell_0}.
\label{eq:Vage}
\end{equation}

Applying this inequality successively gives, for every $n\geq1$,
\begin{equation}
\begin{aligned}
&
\left\|
K(u_0;u_1)\diamond\cdots\diamond
K(u_{n-1};u_n)
\right\|_{-1,-r}
\\
&\qquad\leq
A(r-\ell_0)^{n-1}
\prod_{j=1}^{n}
\|K(u_{j-1};u_j)\|_{-1,-\ell_0}.
\end{aligned}
\label{eq:iterated-vage}
\end{equation}

Combining \eqref{eq:iterated-vage} with
\eqref{eq:K-factor-estimate}, we obtain
\begin{equation}
\begin{aligned}
&
\left\|
K(u_0;u_1)\diamond\cdots\diamond
K(u_{n-1};u_n)
\right\|_{-1,-r}
\\
&\qquad\leq
A(r-\ell_0)^{n-1}
C_{\mathcal K}^{\,n}
\prod_{j=1}^{n}
\frac{(s_{j-1}-s_j)^{\alpha-1}}
{\Gamma(\alpha)}.
\end{aligned}
\label{eq:wick-product-bound}
\end{equation}

Therefore, by the definition of $Y_n$ and the Bochner integral inequality,
\begin{equation}
\begin{aligned}
\|Y_n(t,x)\|_{-1,-r}
\leq{}&
|y_0|
A(r-\ell_0)^{n-1}
C_{\mathcal K}^{\,n}
\\
&\times
\int_{\Delta_n(t,x)}
\prod_{j=1}^{n}
\frac{(s_{j-1}-s_j)^{\alpha-1}}
{\Gamma(\alpha)}
\,du_n\cdots du_1.
\end{aligned}
\label{eq:Yn-bound1}
\end{equation}

The spatial and temporal integrations can be separated. First,
\begin{equation}
\int_{0<z_n<\cdots<z_1<x}
dz_n\cdots dz_1
=
\frac{x^n}{n!}.
\label{eq:spatial-simplex}
\end{equation}

On the other hand, by the iterated fractional integral identity,
\begin{equation}
\begin{aligned}
&
\int_{0<s_n<\cdots<s_1<t}
\prod_{j=1}^{n}
\frac{(s_{j-1}-s_j)^{\alpha-1}}
{\Gamma(\alpha)}
\,ds_n\cdots ds_1
\\
&\qquad=
\frac{t^{n\alpha}}
{\Gamma(n\alpha+1)}.
\end{aligned}
\label{eq:fractional-simplex}
\end{equation}

Consequently,
\begin{equation}
\boxed{
\|Y_n(t,x)\|_{-1,-r}
\leq
|y_0|
A(r-\ell_0)^{n-1}
C_{\mathcal K}^{\,n}
\frac{t^{n\alpha}}
{\Gamma(n\alpha+1)}
\frac{x^n}{n!}.
}
\label{eq:Yn-final-bound}
\end{equation}

Hence,
\begin{equation}
\begin{aligned}
\sup_{(t,x)\in[0,T]\times[0,X]}
\|Y_n(t,x)\|_{-1,-r}
\leq{}&
|y_0|
A(r-\ell_0)^{n-1}
C_{\mathcal K}^{\,n}
\\
&\times
\frac{T^{n\alpha}}
{\Gamma(n\alpha+1)}
\frac{X^n}{n!}.
\end{aligned}
\label{eq:uniform-Yn}
\end{equation}

We now prove the convergence of the numerical majorant. Since
$A(r-\ell_0)$, $C_{\mathcal K}$, $T$ and $X$ are fixed constants, it is enough to observe that
\[
\Gamma(n\alpha+1)n!
\]
grows faster than any exponential sequence. Consequently,
\begin{equation}
\sum_{n=1}^{\infty}
A(r-\ell_0)^{n-1}
C_{\mathcal K}^{\,n}
\frac{T^{n\alpha}X^n}
{\Gamma(n\alpha+1)n!}
<\infty.
\label{eq:majorant-convergence}
\end{equation}

It follows from \eqref{eq:uniform-Yn} and the Weierstrass
$M$-test that
\begin{equation}
\sum_{n=1}^{\infty}Y_n(t,x)
\end{equation}
converges uniformly on $[0,T]\times[0,X]$ in the norm of
$(\mathcal S)_{-1,-r}$.

It remains to prove continuity.

For every fixed $n$, the map
\[
(t,x)\longmapsto Y_n(t,x)
\]
is continuous with values in $(\mathcal S)_{-1,-r}$.

Indeed, by assumption,
\[
(s,z)\longmapsto\mathcal K(s,z)
\]
is continuous with values in $(\mathcal S)_{-1,-\ell_0}$.
The Wick product is continuous under the V$\mathring{a}$ge estimate, and the
integrand defining $Y_n$ is dominated by the integrable function
appearing on the right-hand side of \eqref{eq:wick-product-bound}.
Since $0<\alpha<1$, the singularity
$(s_{j-1}-s_j)^{\alpha-1}$ is integrable on the ordered simplex.
Therefore, the dominated convergence theorem for Bochner integrals
implies that
\[
Y_n\in
C\big([0,T]\times[0,X];(\mathcal S)_{-1,-r}\big).
\]

Since the series $\sum_{n\geq1}Y_n$ converges uniformly and each
$Y_n$ is continuous, its sum is continuous. Thus
\begin{equation}
\boxed{
Y\in
C\big([0,T]\times[0,X];(\mathcal S)_{-1,-r}\big).
}
\label{eq:Y-continuity}
\end{equation}

Because $[0,T]\times[0,X]$ is compact and $Y$ is continuous,
we immediately obtain
\begin{equation}
\boxed{
M_Y:=
\sup_{(t,x)\in[0,T]\times[0,X]}
\|Y(t,x)\|_{-1,-r}
<\infty.
}
\label{eq:Y-bound}
\end{equation}

We emphasize that \eqref{eq:Y-bound} is not an additional assumption.
It follows from the already established convergence of the series in
the space
$C([0,T]\times[0,X];(\mathcal S)_{-1,-r})$.

We now estimate the remainder. Iterating \eqref{eq:Volterra} $n$
times yields
\begin{equation}
Y(t,x)
=
y_0+\sum_{j=1}^{n}Y_j(t,x)+R_n(t,x),
\label{eq:remainder-decomposition}
\end{equation}
where
\begin{equation}
\begin{aligned}
R_n(t,x)
={}&
\int_{\Delta_{n+1}(t,x)}
K(u_0;u_1)\diamond\cdots\diamond
K(u_n;u_{n+1})
\\
&\qquad\diamond
Y(u_{n+1})
\,du_{n+1}\cdots du_1.
\end{aligned}
\label{eq:Rn-definition}
\end{equation}

There are $n+1$ kernel factors in \eqref{eq:Rn-definition}.
Since
\[
Y(u_{n+1})\in(\mathcal S)_{-1,-r}
\]
and
\[
\|Y(u_{n+1})\|_{-1,-r}\leq M_Y,
\]
successive applications of the V\ aage inequality yield
\begin{equation}
\begin{aligned}
&
\left\|
K(u_0;u_1)\diamond\cdots\diamond
K(u_n;u_{n+1})
\diamond Y(u_{n+1})
\right\|_{-1,-r}
\\
&\qquad\leq
M_Y
A(r-\ell_0)^{n+1}
C_{\mathcal K}^{\,n+1}
\prod_{j=1}^{n+1}
\frac{(s_{j-1}-s_j)^{\alpha-1}}
{\Gamma(\alpha)}.
\end{aligned}
\label{eq:Rn-integrand-bound}
\end{equation}

Consequently,
\begin{equation}
\begin{aligned}
\|R_n(t,x)\|_{-1,-r}
\leq{}&
M_Y
A(r-\ell_0)^{n+1}
C_{\mathcal K}^{\,n+1}
\\
&\times
\int_{\Delta_{n+1}(t,x)}
\prod_{j=1}^{n+1}
\frac{(s_{j-1}-s_j)^{\alpha-1}}
{\Gamma(\alpha)}
\,du_{n+1}\cdots du_1.
\end{aligned}
\label{eq:Rn-bound1}
\end{equation}

Using again the identities
\begin{equation}
\int_{0<z_{n+1}<\cdots<z_1<x}
dz_{n+1}\cdots dz_1
=
\frac{x^{n+1}}{(n+1)!}
\end{equation}
and
\begin{equation}
\begin{aligned}
&
\int_{0<s_{n+1}<\cdots<s_1<t}
\prod_{j=1}^{n+1}
\frac{(s_{j-1}-s_j)^{\alpha-1}}
{\Gamma(\alpha)}
\,ds_{n+1}\cdots ds_1
\\
&\qquad=
\frac{t^{(n+1)\alpha}}
{\Gamma((n+1)\alpha+1)},
\end{aligned}
\end{equation}
we obtain
\begin{equation}
\boxed{
\begin{aligned}
\|R_n(t,x)\|_{-1,-r}
\leq{}&
M_Y
A(r-\ell_0)^{n+1}
C_{\mathcal K}^{\,n+1}
\\
&\times
\frac{t^{(n+1)\alpha}}
{\Gamma((n+1)\alpha+1)}
\frac{x^{n+1}}
{(n+1)!}.
\end{aligned}
}
\label{eq:Rn-final}
\end{equation}

Therefore,
\begin{equation}
\begin{aligned}
\sup_{(t,x)\in[0,T]\times[0,X]}
\|R_n(t,x)\|_{-1,-r}
\leq{}&
M_Y
A(r-\ell_0)^{n+1}
C_{\mathcal K}^{\,n+1}
\\
&\times
\frac{T^{(n+1)\alpha}}
{\Gamma((n+1)\alpha+1)}
\frac{X^{n+1}}
{(n+1)!}.
\end{aligned}
\label{eq:Rn-uniform}
\end{equation}

The right-hand side of \eqref{eq:Rn-uniform} converges to zero as
$n\to\infty$. Hence
\begin{equation}
\boxed{
\sup_{(t,x)\in[0,T]\times[0,X]}
\|R_n(t,x)\|_{-1,-r}
\longrightarrow0.
}
\label{eq:Rn-to-zero}
\end{equation}

Thus the infinite iterated expansion is valid in
$(\mathcal S)_{-1,-r}$ and satisfies the original Volterra equation.

In particular,
\[
Y(t,x)
=
y_0+\sum_{n=1}^{\infty}Y_n(t,x)
\]
converges uniformly on $[0,T]\times[0,X]$, and
\[
Y\in
C\big([0,T]\times[0,X];(\mathcal S)_{-1,-r}\big).
\]

Finally, we stress that no smallness condition of the form
\[
A(r-\ell_0)C_{\mathcal K}<1
\]
is required. The convergence follows from the Gamma-function and
factorial terms
\[
\frac{T^{n\alpha}}{\Gamma(n\alpha+1)}
\qquad\text{and}\qquad
\frac{X^n}{n!},
\]
which dominate the exponential growth produced by the successive
applications of the V$\mathring{a}$ge inequality.
\end{proof}
```


